%%%%%%%%%%%%%%%%%%%%%%%%%%%%%%%%%%%%%%%%%%%%%%%%%%%%%%%%%%%%
% 3.8 ALGEBRAIC GEOMETRY
% The Composition of Advanced Mathematics
%%%%%%%%%%%%%%%%%%%%%%%%%%%%%%%%%%%%%%%%%%%%%%%%%%%%%%%%%%%%

\section{Algebraic Geometry}
\label{sec:algebraic-geometry}

\prerequisite{
Elementary Algebra from \cref{sec:elementary-algebra},
Abstract Algebra from \cref{sec:abstract-algebra},
Ring Theory from \cref{sec:ring-theory},
Field Theory from \cref{sec:field-theory},
Commutative Algebra from \cref{sec:commutative-algebra},
and Projective Geometry from \cref{sec:projective-geometry}.
}

\learningobjectives{
By the end of this section, the reader should be able to:
\begin{itemize}
    \item explain the central goals of algebraic geometry;
    \item interpret polynomial equations as geometric objects;
    \item define affine algebraic sets and affine varieties;
    \item compute basic vanishing sets \(V(I)\);
    \item determine ideals of algebraic sets;
    \item explain the algebra--geometry correspondence;
    \item state the geometric meaning of Hilbert's Nullstellensatz;
    \item describe the Zariski topology;
    \item distinguish reducible and irreducible algebraic sets;
    \item work with coordinate rings;
    \item identify singular and nonsingular points;
    \item explain tangent spaces using polynomial derivatives;
    \item describe projective algebraic sets using homogeneous polynomials;
    \item explain why projective space provides points at infinity;
    \item recognize rational maps and birational equivalence;
    \item distinguish curves, surfaces, and higher-dimensional varieties;
    \item explain the motivation for schemes and the spectrum of a ring;
    \item recognize the roles of sheaves and local rings;
    \item connect algebraic geometry with number theory, topology,
    complex geometry, and modern mathematics.
\end{itemize}
}

%%%%%%%%%%%%%%%%%%%%%%%%%%%%%%%%%%%%%%%%%%%%%%%%%%%%%%%%%%%%
\subsection{What Is Algebraic Geometry?}
\label{subsec:what-is-algebraic-geometry}
%%%%%%%%%%%%%%%%%%%%%%%%%%%%%%%%%%%%%%%%%%%%%%%%%%%%%%%%%%%%

Algebraic geometry studies geometric objects defined by polynomial
equations.

A simple example is the circle

\[
x^2+y^2=1.
\]

A more complicated example is the cubic curve

\[
y^2=x^3-x.
\]

Rather than treating these merely as equations to solve, algebraic geometry
treats their complete solution sets as geometric spaces.

\begin{definitionbox}{Algebraic Geometry}
\emph{Algebraic geometry} is the study of geometric spaces defined by
polynomial equations and the algebraic structures associated with those
spaces.
\end{definitionbox}

Its central principle can be summarized as

\[
\boxed{
\text{polynomial algebra}
\quad\longleftrightarrow\quad
\text{geometry}.
}
\]

%%%%%%%%%%%%%%%%%%%%%%%%%%%%%%%%%%%%%%%%%%%%%%%%%%%%%%%%%%%%
\subsection{Polynomial Equations as Geometry}
\label{subsec:polynomial-equations-as-geometry}
%%%%%%%%%%%%%%%%%%%%%%%%%%%%%%%%%%%%%%%%%%%%%%%%%%%%%%%%%%%%

Consider a polynomial

\[
f(x,y)=x^2+y^2-1.
\]

The equation

\[
f(x,y)=0
\]

defines a geometric set.

Over the real numbers,

\[
V(f)
=
\left\{
(x,y)\in\R^2:
x^2+y^2-1=0
\right\},
\]

which is the unit circle.

However, algebraic geometry is not restricted to \(\R\).

The same polynomial may be studied over

\[
\C,
\qquad
\Q,
\qquad
\F_p,
\]

or more general fields.

The resulting geometry depends strongly on the underlying field.

%%%%%%%%%%%%%%%%%%%%%%%%%%%%%%%%%%%%%%%%%%%%%%%%%%%%%%%%%%%%
\subsection{Affine Space}
\label{subsec:algebraic-affine-space}
%%%%%%%%%%%%%%%%%%%%%%%%%%%%%%%%%%%%%%%%%%%%%%%%%%%%%%%%%%%%

Let \(k\) be a field.

\begin{definitionbox}{Affine Space}
The \(n\)-dimensional affine space over \(k\) is

\[
\boxed{
\A^n_k=k^n.
}
\]

Its points are ordered \(n\)-tuples

\[
(a_1,\ldots,a_n)
\]

with

\[
a_i\in k.
\]
\end{definitionbox}

The symbol

\[
\A^n_k
\]

is read ``affine \(n\)-space over \(k\).''

Examples include

\[
\A^1_{\R}=\R,
\]

\[
\A^2_{\R}=\R^2,
\]

and

\[
\A^n_{\C}=\C^n.
\]

%%%%%%%%%%%%%%%%%%%%%%%%%%%%%%%%%%%%%%%%%%%%%%%%%%%%%%%%%%%%
\subsection{Polynomial Rings}
\label{subsec:geometry-polynomial-rings}
%%%%%%%%%%%%%%%%%%%%%%%%%%%%%%%%%%%%%%%%%%%%%%%%%%%%%%%%%%%%

Polynomial equations in \(n\) variables over a field \(k\) belong to the
polynomial ring

\[
\boxed{
k[x_1,\ldots,x_n].
}
\]

This ring becomes the algebraic language for studying subsets of

\[
\A^n_k.
\]

For example,

\[
x^2+y^2-1
\]

belongs to

\[
\R[x,y].
\]

A system such as

\[
x^2+y^2-1=0,
\]

\[
x-y=0
\]

corresponds to the pair of polynomials

\[
x^2+y^2-1,
\qquad
x-y.
\]

%%%%%%%%%%%%%%%%%%%%%%%%%%%%%%%%%%%%%%%%%%%%%%%%%%%%%%%%%%%%
\subsection{Vanishing Sets}
\label{subsec:vanishing-sets}
%%%%%%%%%%%%%%%%%%%%%%%%%%%%%%%%%%%%%%%%%%%%%%%%%%%%%%%%%%%%

Let

\[
S\subseteq k[x_1,\ldots,x_n]
\]

be a collection of polynomials.

\begin{definitionbox}{Vanishing Set}
The \emph{vanishing set} of \(S\) is

\[
\boxed{
V(S)
=
\left\{
p\in\A^n_k:
f(p)=0
\text{ for every }f\in S
\right\}.
}
\]
\end{definitionbox}

The notation

\[
V(S)
\]

is read ``the vanishing set of \(S\)'' or ``V of \(S\).''

Thus \(V(S)\) consists of all common solutions to the polynomial equations
in \(S\).

%%%%%%%%%%%%%%%%%%%%%%%%%%%%%%%%%%%%%%%%%%%%%%%%%%%%%%%%%%%%
\subsection{Examples of Vanishing Sets}
\label{subsec:vanishing-set-examples}
%%%%%%%%%%%%%%%%%%%%%%%%%%%%%%%%%%%%%%%%%%%%%%%%%%%%%%%%%%%%

For

\[
f(x,y)=y-x^2,
\]

we obtain

\[
\boxed{
V(f)=\{(x,y):y=x^2\},
}
\]

which is a parabola over \(\R\).

For

\[
f(x,y)=xy,
\]

we obtain

\[
\boxed{
V(xy)=V(x)\cup V(y),
}
\]

which consists of the two coordinate axes.

For

\[
f(x,y)=x^2+y^2+1
\]

over \(\R\),

\[
\boxed{
V(f)=\varnothing.
}
\]

Over \(\C\), however, the same polynomial has solutions.

%%%%%%%%%%%%%%%%%%%%%%%%%%%%%%%%%%%%%%%%%%%%%%%%%%%%%%%%%%%%
\subsection{Systems of Polynomial Equations}
\label{subsec:systems-polynomial-equations}
%%%%%%%%%%%%%%%%%%%%%%%%%%%%%%%%%%%%%%%%%%%%%%%%%%%%%%%%%%%%

Suppose

\[
S=
\{
f_1,\ldots,f_m
\}.
\]

Then

\[
\boxed{
V(S)
=
V(f_1,\ldots,f_m)
}
\]

is the common solution set

\[
f_1=0,
\ldots,
f_m=0.
\]

For example,

\[
V(x-y,x+y-2)
\]

consists of points satisfying

\[
x=y
\]

and

\[
x+y=2.
\]

Solving gives

\[
x=y=1.
\]

Therefore

\[
\boxed{
V(x-y,x+y-2)=\{(1,1)\}.
}
\]

%%%%%%%%%%%%%%%%%%%%%%%%%%%%%%%%%%%%%%%%%%%%%%%%%%%%%%%%%%%%
\subsection{Why Ideals Appear}
\label{subsec:why-ideals-appear}
%%%%%%%%%%%%%%%%%%%%%%%%%%%%%%%%%%%%%%%%%%%%%%%%%%%%%%%%%%%%

Suppose a point \(p\) satisfies

\[
f(p)=0
\]

and

\[
g(p)=0.
\]

Then it also satisfies

\[
(f+g)(p)=0.
\]

Furthermore, for any polynomial \(h\),

\[
(hf)(p)=0.
\]

Therefore the polynomials vanishing on a common solution set naturally form
an ideal.

If

\[
I=(f_1,\ldots,f_m),
\]

then

\[
\boxed{
V(f_1,\ldots,f_m)=V(I).
}
\]

Thus ideals, rather than arbitrary lists of equations, are the natural
algebraic objects associated with algebraic sets.

%%%%%%%%%%%%%%%%%%%%%%%%%%%%%%%%%%%%%%%%%%%%%%%%%%%%%%%%%%%%
\subsection{Affine Algebraic Sets}
\label{subsec:affine-algebraic-sets}
%%%%%%%%%%%%%%%%%%%%%%%%%%%%%%%%%%%%%%%%%%%%%%%%%%%%%%%%%%%%

\begin{definitionbox}{Affine Algebraic Set}
A subset

\[
X\subseteq\A^n_k
\]

is an \emph{affine algebraic set} if there exists an ideal

\[
I\subseteq k[x_1,\ldots,x_n]
\]

such that

\[
\boxed{
X=V(I).
}
\]
\end{definitionbox}

Thus affine algebraic sets are precisely the common zero sets of collections
of polynomials.

%%%%%%%%%%%%%%%%%%%%%%%%%%%%%%%%%%%%%%%%%%%%%%%%%%%%%%%%%%%%
\subsection{Geometric Operations and Ideals}
\label{subsec:geometric-operations-ideals}
%%%%%%%%%%%%%%%%%%%%%%%%%%%%%%%%%%%%%%%%%%%%%%%%%%%%%%%%%%%%

The algebra of ideals translates into geometry.

For ideals \(I,J\),

\[
\boxed{
V(I+J)=V(I)\cap V(J).
}
\]

Also,

\[
\boxed{
V(IJ)=V(I)\cup V(J).
}
\]

Similarly,

\[
\boxed{
V(I\cap J)=V(I)\cup V(J).
}
\]

This reversal is fundamental:

\[
\boxed{
\text{larger ideals}
\quad\longleftrightarrow\quad
\text{smaller algebraic sets}.
}
\]

%%%%%%%%%%%%%%%%%%%%%%%%%%%%%%%%%%%%%%%%%%%%%%%%%%%%%%%%%%%%
\subsection{Worked Example: Intersecting a Line and Parabola}
\label{subsec:worked-line-parabola}
%%%%%%%%%%%%%%%%%%%%%%%%%%%%%%%%%%%%%%%%%%%%%%%%%%%%%%%%%%%%

\begin{workedexamplebox}{Intersection of Two Algebraic Sets}

Consider

\[
V(y-x^2)
\]

and

\[
V(y-x).
\]

Their intersection satisfies

\[
y=x^2
\]

and

\[
y=x.
\]

Therefore

\[
x^2=x,
\]

so

\[
x(x-1)=0.
\]

Hence

\[
x=0
\qquad\text{or}\qquad
x=1.
\]

The corresponding points are

\[
(0,0)
\]

and

\[
(1,1).
\]

\begin{answerbox}
\[
\boxed{
V(y-x^2)\cap V(y-x)
=
\{(0,0),(1,1)\}
}
\]
\end{answerbox}
\end{workedexamplebox}

%%%%%%%%%%%%%%%%%%%%%%%%%%%%%%%%%%%%%%%%%%%%%%%%%%%%%%%%%%%%
\subsection{Ideals of Algebraic Sets}
\label{subsec:ideals-of-algebraic-sets}
%%%%%%%%%%%%%%%%%%%%%%%%%%%%%%%%%%%%%%%%%%%%%%%%%%%%%%%%%%%%

The correspondence can be reversed.

Given

\[
X\subseteq\A^n_k,
\]

consider every polynomial that vanishes on \(X\).

\begin{definitionbox}{Vanishing Ideal}
The \emph{vanishing ideal} of \(X\) is

\[
\boxed{
I(X)
=
\left\{
f\in k[x_1,\ldots,x_n]:
f(p)=0
\text{ for every }p\in X
\right\}.
}
\]
\end{definitionbox}

The notation

\[
I(X)
\]

is read ``the ideal of \(X\).''

%%%%%%%%%%%%%%%%%%%%%%%%%%%%%%%%%%%%%%%%%%%%%%%%%%%%%%%%%%%%
\subsection{The Algebra--Geometry Dictionary}
\label{subsec:algebra-geometry-dictionary}
%%%%%%%%%%%%%%%%%%%%%%%%%%%%%%%%%%%%%%%%%%%%%%%%%%%%%%%%%%%%

The two basic operations are

\[
\boxed{
I
\longmapsto
V(I)
}
\]

and

\[
\boxed{
X
\longmapsto
I(X).
}
\]

They reverse inclusions.

If

\[
I\subseteq J,
\]

then

\[
\boxed{
V(J)\subseteq V(I).
}
\]

If

\[
X\subseteq Y,
\]

then

\[
\boxed{
I(Y)\subseteq I(X).
}
\]

This contravariant behavior is one of the organizing principles of
algebraic geometry.

%%%%%%%%%%%%%%%%%%%%%%%%%%%%%%%%%%%%%%%%%%%%%%%%%%%%%%%%%%%%
\subsection{Radical Ideals}
\label{subsec:geometry-radical-ideals}
%%%%%%%%%%%%%%%%%%%%%%%%%%%%%%%%%%%%%%%%%%%%%%%%%%%%%%%%%%%%

Suppose

\[
f^m
\]

vanishes at a point.

Then

\[
f
\]

also vanishes there.

Therefore ordinary zero sets cannot distinguish between equations such as

\[
f=0
\]

and

\[
f^m=0.
\]

This motivates radical ideals.

Recall that

\[
\boxed{
\sqrt{I}
=
\{
f:
f^m\in I
\text{ for some }m\geq1
\}.
}
\]

The ideal

\[
\sqrt{I}
\]

is called the \emph{radical} of \(I\).

Moreover,

\[
\boxed{
V(I)=V(\sqrt{I}).
}
\]

%%%%%%%%%%%%%%%%%%%%%%%%%%%%%%%%%%%%%%%%%%%%%%%%%%%%%%%%%%%%
\subsection{Example: Same Geometry, Different Ideals}
\label{subsec:same-geometry-different-ideals}
%%%%%%%%%%%%%%%%%%%%%%%%%%%%%%%%%%%%%%%%%%%%%%%%%%%%%%%%%%%%

Consider

\[
I=(x)
\]

and

\[
J=(x^2)
\]

in

\[
k[x,y].
\]

Both have the same ordinary vanishing set:

\[
\boxed{
V(x)=V(x^2).
}
\]

Geometrically, both describe the \(y\)-axis.

But algebraically,

\[
(x)\neq(x^2).
\]

This loss of information is one motivation for the later scheme-theoretic
viewpoint.

%%%%%%%%%%%%%%%%%%%%%%%%%%%%%%%%%%%%%%%%%%%%%%%%%%%%%%%%%%%%
\subsection{Hilbert's Nullstellensatz}
\label{subsec:nullstellensatz}
%%%%%%%%%%%%%%%%%%%%%%%%%%%%%%%%%%%%%%%%%%%%%%%%%%%%%%%%%%%%

Over an algebraically closed field, the relationship between ideals and
algebraic sets becomes especially powerful.

\begin{theorem}[Hilbert's Nullstellensatz --- Geometric Form]
Let \(k\) be algebraically closed and let

\[
I\subseteq k[x_1,\ldots,x_n].
\]

Then

\[
\boxed{
I(V(I))=\sqrt{I}.
}
\]
\end{theorem}

The German word \emph{Nullstellensatz} means approximately ``zero-locus
theorem.''

The theorem states that the polynomials vanishing on the solution set of
\(I\) are exactly the elements of the radical of \(I\).

%%%%%%%%%%%%%%%%%%%%%%%%%%%%%%%%%%%%%%%%%%%%%%%%%%%%%%%%%%%%
\subsection{Weak Nullstellensatz}
\label{subsec:weak-nullstellensatz}
%%%%%%%%%%%%%%%%%%%%%%%%%%%%%%%%%%%%%%%%%%%%%%%%%%%%%%%%%%%%

A closely related form states:

\begin{theorem}[Weak Nullstellensatz]
If \(k\) is algebraically closed, then every maximal ideal of

\[
k[x_1,\ldots,x_n]
\]

has the form

\[
\boxed{
(x_1-a_1,\ldots,x_n-a_n)
}
\]

for some

\[
(a_1,\ldots,a_n)\in k^n.
\]
\end{theorem}

Thus points of affine space correspond to maximal ideals.

This is an important early appearance of the principle

\[
\boxed{
\text{geometric points}
\quad\longleftrightarrow\quad
\text{algebraic ideals}.
}
\]

%%%%%%%%%%%%%%%%%%%%%%%%%%%%%%%%%%%%%%%%%%%%%%%%%%%%%%%%%%%%
\subsection{The Zariski Topology}
\label{subsec:zariski-topology}
%%%%%%%%%%%%%%%%%%%%%%%%%%%%%%%%%%%%%%%%%%%%%%%%%%%%%%%%%%%%

Algebraic sets can be used as the closed sets of a topology.

\begin{definitionbox}{Zariski Topology}
The \emph{Zariski topology} on \(\A^n_k\) is the topology whose closed sets
are the affine algebraic sets

\[
\boxed{
V(I).
}
\]
\end{definitionbox}

This topology is very different from the ordinary Euclidean topology.

%%%%%%%%%%%%%%%%%%%%%%%%%%%%%%%%%%%%%%%%%%%%%%%%%%%%%%%%%%%%
\subsection{Why the Zariski Topology Is Coarse}
\label{subsec:zariski-coarse}
%%%%%%%%%%%%%%%%%%%%%%%%%%%%%%%%%%%%%%%%%%%%%%%%%%%%%%%%%%%%

Consider

\[
\A^1_k.
\]

If \(k\) is an infinite field, a nonzero one-variable polynomial has only
finitely many roots.

Therefore the proper algebraic subsets of

\[
\A^1_k
\]

are finite sets.

Consequently, over an algebraically closed infinite field, the closed sets
of the affine line are

\[
\boxed{
\varnothing,
\quad
\text{finite subsets},
\quad
\A^1_k.
}
\]

Thus Zariski-open sets are generally very large.

%%%%%%%%%%%%%%%%%%%%%%%%%%%%%%%%%%%%%%%%%%%%%%%%%%%%%%%%%%%%
\subsection{Basic Open Sets}
\label{subsec:basic-open-sets}
%%%%%%%%%%%%%%%%%%%%%%%%%%%%%%%%%%%%%%%%%%%%%%%%%%%%%%%%%%%%

For a polynomial

\[
f\in k[x_1,\ldots,x_n],
\]

define

\[
\boxed{
D(f)
=
\A^n_k\setminus V(f).
}
\]

The notation

\[
D(f)
\]

is read ``D of \(f\).''

Thus

\[
D(f)
=
\{
p:f(p)\neq0
\}.
\]

These sets form a basis for the Zariski topology on affine space.

%%%%%%%%%%%%%%%%%%%%%%%%%%%%%%%%%%%%%%%%%%%%%%%%%%%%%%%%%%%%
\subsection{Irreducibility}
\label{subsec:algebraic-irreducibility}
%%%%%%%%%%%%%%%%%%%%%%%%%%%%%%%%%%%%%%%%%%%%%%%%%%%%%%%%%%%%

Some algebraic sets can be decomposed into simpler algebraic pieces.

Others cannot.

\begin{definitionbox}{Irreducible Algebraic Set}
A nonempty algebraic set \(X\) is \emph{irreducible} if it cannot be written

\[
\boxed{
X=X_1\cup X_2
}
\]

where \(X_1\) and \(X_2\) are proper closed subsets of \(X\).
\end{definitionbox}

An irreducible algebraic set is traditionally called an
\emph{affine variety} in one common convention.

Terminology varies among texts.

%%%%%%%%%%%%%%%%%%%%%%%%%%%%%%%%%%%%%%%%%%%%%%%%%%%%%%%%%%%%
\subsection{Reducible Example}
\label{subsec:reducible-example}
%%%%%%%%%%%%%%%%%%%%%%%%%%%%%%%%%%%%%%%%%%%%%%%%%%%%%%%%%%%%

Consider

\[
X=V(xy).
\]

Since

\[
xy=0
\]

means

\[
x=0
\qquad\text{or}\qquad
y=0,
\]

we obtain

\[
\boxed{
V(xy)=V(x)\cup V(y).
}
\]

Thus \(X\) is reducible.

Its irreducible components are the two coordinate axes.

%%%%%%%%%%%%%%%%%%%%%%%%%%%%%%%%%%%%%%%%%%%%%%%%%%%%%%%%%%%%
\subsection{Prime Ideals and Irreducible Geometry}
\label{subsec:prime-ideals-irreducible}
%%%%%%%%%%%%%%%%%%%%%%%%%%%%%%%%%%%%%%%%%%%%%%%%%%%%%%%%%%%%

Over an algebraically closed field, irreducibility has a precise algebraic
translation.

\begin{theorem}
An affine algebraic set \(X\) is irreducible if and only if its vanishing
ideal

\[
I(X)
\]

is a prime ideal.
\end{theorem}

Thus

\[
\boxed{
\text{irreducible geometry}
\quad\longleftrightarrow\quad
\text{prime ideals}.
}
\]

This connection becomes even more important in scheme theory.

%%%%%%%%%%%%%%%%%%%%%%%%%%%%%%%%%%%%%%%%%%%%%%%%%%%%%%%%%%%%
\subsection{Coordinate Rings}
\label{subsec:coordinate-rings}
%%%%%%%%%%%%%%%%%%%%%%%%%%%%%%%%%%%%%%%%%%%%%%%%%%%%%%%%%%%%

Suppose

\[
X=V(I)\subseteq\A^n_k.
\]

\begin{definitionbox}{Coordinate Ring}
The \emph{coordinate ring} of \(X\) is

\[
\boxed{
k[X]
=
k[x_1,\ldots,x_n]/I(X).
}
\]
\end{definitionbox}

The coordinate ring consists of polynomial functions on \(X\), where two
polynomials are identified when they define the same function on \(X\).

%%%%%%%%%%%%%%%%%%%%%%%%%%%%%%%%%%%%%%%%%%%%%%%%%%%%%%%%%%%%
\subsection{Example: Coordinate Ring of a Parabola}
\label{subsec:coordinate-ring-parabola}
%%%%%%%%%%%%%%%%%%%%%%%%%%%%%%%%%%%%%%%%%%%%%%%%%%%%%%%%%%%%

Let

\[
X=V(y-x^2)\subseteq\A^2_k.
\]

Then

\[
\boxed{
k[X]
=
k[x,y]/(y-x^2).
}
\]

Inside the quotient,

\[
y=x^2.
\]

Therefore every polynomial expression involving \(x\) and \(y\) can be
reduced to a polynomial involving \(x\) alone.

In fact,

\[
\boxed{
k[X]\cong k[x].
}
\]

Algebraically, the parabola behaves like the affine line.

%%%%%%%%%%%%%%%%%%%%%%%%%%%%%%%%%%%%%%%%%%%%%%%%%%%%%%%%%%%%
\subsection{Regular Functions}
\label{subsec:regular-functions}
%%%%%%%%%%%%%%%%%%%%%%%%%%%%%%%%%%%%%%%%%%%%%%%%%%%%%%%%%%%%

Polynomial functions on affine varieties are examples of
\emph{regular functions}.

On an affine algebraic set \(X\), a regular function is locally expressible
as

\[
\boxed{
\frac{f}{g}
}
\]

where \(f\) and \(g\) are polynomials and \(g\) does not vanish on the
relevant neighborhood.

For affine varieties, global regular functions are represented by elements
of the coordinate ring under the standard hypotheses.

%%%%%%%%%%%%%%%%%%%%%%%%%%%%%%%%%%%%%%%%%%%%%%%%%%%%%%%%%%%%
\subsection{Morphisms of Affine Varieties}
\label{subsec:morphisms-affine-varieties}
%%%%%%%%%%%%%%%%%%%%%%%%%%%%%%%%%%%%%%%%%%%%%%%%%%%%%%%%%%%%

A polynomial map

\[
F:\A^n_k\to\A^m_k
\]

has the form

\[
\boxed{
F(p)
=
\bigl(
f_1(p),\ldots,f_m(p)
\bigr),
}
\]

where

\[
f_i\in k[x_1,\ldots,x_n].
\]

When such a map sends one algebraic variety into another, it gives a basic
example of a morphism.

Morphisms are the structure-preserving maps of algebraic geometry.

%%%%%%%%%%%%%%%%%%%%%%%%%%%%%%%%%%%%%%%%%%%%%%%%%%%%%%%%%%%%
\subsection{Contravariance of Coordinate Rings}
\label{subsec:contravariance-coordinate-rings}
%%%%%%%%%%%%%%%%%%%%%%%%%%%%%%%%%%%%%%%%%%%%%%%%%%%%%%%%%%%%

A morphism

\[
F:X\to Y
\]

induces a map in the opposite direction on functions:

\[
\boxed{
F^*:k[Y]\to k[X].
}
\]

It is defined by

\[
\boxed{
F^*(f)=f\circ F.
}
\]

The map \(F^*\) is called the \emph{pullback}.

Thus geometry and algebra again reverse direction:

\[
\boxed{
X\longrightarrow Y
\qquad\Longleftrightarrow\qquad
k[Y]\longrightarrow k[X].
}
\]

%%%%%%%%%%%%%%%%%%%%%%%%%%%%%%%%%%%%%%%%%%%%%%%%%%%%%%%%%%%%
\subsection{Dimension}
\label{subsec:algebraic-dimension}
%%%%%%%%%%%%%%%%%%%%%%%%%%%%%%%%%%%%%%%%%%%%%%%%%%%%%%%%%%%%

Dimension measures the number of independent parameters needed to describe a
geometric object locally or generically.

Typical examples are

\[
\boxed{
\dim(\text{point})=0,
}
\]

\[
\boxed{
\dim(\text{curve})=1,
}
\]

\[
\boxed{
\dim(\text{surface})=2.
}
\]

In algebraic geometry, dimension can be defined algebraically using chains of
prime ideals.

%%%%%%%%%%%%%%%%%%%%%%%%%%%%%%%%%%%%%%%%%%%%%%%%%%%%%%%%%%%%
\subsection{Krull Dimension}
\label{subsec:krull-dimension-geometry}
%%%%%%%%%%%%%%%%%%%%%%%%%%%%%%%%%%%%%%%%%%%%%%%%%%%%%%%%%%%%

For a commutative ring \(R\), the Krull dimension is the supremum of lengths
of chains

\[
\boxed{
\mathfrak p_0
\subsetneq
\mathfrak p_1
\subsetneq
\cdots
\subsetneq
\mathfrak p_r
}
\]

of prime ideals.

For an affine variety \(X\),

\[
\boxed{
\dim X=\dim k[X].
}
\]

This is a remarkable example of a geometric quantity being encoded entirely
in a ring.

%%%%%%%%%%%%%%%%%%%%%%%%%%%%%%%%%%%%%%%%%%%%%%%%%%%%%%%%%%%%
\subsection{Hypersurfaces}
\label{subsec:algebraic-hypersurfaces}
%%%%%%%%%%%%%%%%%%%%%%%%%%%%%%%%%%%%%%%%%%%%%%%%%%%%%%%%%%%%

\begin{definitionbox}{Hypersurface}
A \emph{hypersurface} in \(\A^n_k\) is an algebraic set defined by a single
nonconstant polynomial equation

\[
\boxed{
f(x_1,\ldots,x_n)=0.
}
\]
\end{definitionbox}

Under suitable nondegeneracy assumptions, a hypersurface has dimension

\[
\boxed{
n-1.
}
\]

Examples include plane algebraic curves in \(\A^2\) and algebraic surfaces
in \(\A^3\).

%%%%%%%%%%%%%%%%%%%%%%%%%%%%%%%%%%%%%%%%%%%%%%%%%%%%%%%%%%%%
\subsection{Algebraic Curves}
\label{subsec:algebraic-curves}
%%%%%%%%%%%%%%%%%%%%%%%%%%%%%%%%%%%%%%%%%%%%%%%%%%%%%%%%%%%%

An algebraic curve is, roughly, a one-dimensional algebraic variety.

Plane curves often arise from equations

\[
\boxed{
f(x,y)=0.
}
\]

Examples include

\[
y-x^2=0,
\]

\[
x^2+y^2-1=0,
\]

and

\[
y^2-x^3+x=0.
\]

The geometry of algebraic curves depends on degree, singularities, genus,
and the underlying field.

%%%%%%%%%%%%%%%%%%%%%%%%%%%%%%%%%%%%%%%%%%%%%%%%%%%%%%%%%%%%
\subsection{Algebraic Surfaces}
\label{subsec:algebraic-surfaces}
%%%%%%%%%%%%%%%%%%%%%%%%%%%%%%%%%%%%%%%%%%%%%%%%%%%%%%%%%%%%

An algebraic surface is a two-dimensional algebraic variety.

A hypersurface in affine three-space may be given by

\[
\boxed{
f(x,y,z)=0.
}
\]

Examples include

\[
x^2+y^2+z^2-1=0
\]

and

\[
z-xy=0.
\]

Over \(\R\), these equations can be visualized as familiar surfaces.

Over \(\C\), their real dimension is typically twice their complex
dimension, so ordinary three-dimensional visualization becomes inadequate.

%%%%%%%%%%%%%%%%%%%%%%%%%%%%%%%%%%%%%%%%%%%%%%%%%%%%%%%%%%%%
\subsection{Singular Points}
\label{subsec:algebraic-singular-points}
%%%%%%%%%%%%%%%%%%%%%%%%%%%%%%%%%%%%%%%%%%%%%%%%%%%%%%%%%%%%

Not every algebraic variety is smooth.

Consider a hypersurface

\[
X=V(f).
\]

A point \(p\in X\) is singular when the gradient vanishes there:

\[
\boxed{
\nabla f(p)=\mathbf{0}.
}
\]

Equivalently,

\[
\boxed{
\frac{\partial f}{\partial x_1}(p)
=
\cdots
=
\frac{\partial f}{\partial x_n}(p)
=
0.
}
\]

A point where the gradient is nonzero is nonsingular or smooth.

%%%%%%%%%%%%%%%%%%%%%%%%%%%%%%%%%%%%%%%%%%%%%%%%%%%%%%%%%%%%
\subsection{Example: A Singular Curve}
\label{subsec:singular-curve-example}
%%%%%%%%%%%%%%%%%%%%%%%%%%%%%%%%%%%%%%%%%%%%%%%%%%%%%%%%%%%%

Consider

\[
f(x,y)=y^2-x^3.
\]

Then

\[
f_x=-3x^2,
\]

and

\[
f_y=2y.
\]

At

\[
(0,0),
\]

we have

\[
f_x(0,0)=0,
\qquad
f_y(0,0)=0.
\]

Therefore

\[
\boxed{
(0,0)
\text{ is singular.}
}
\]

The curve

\[
y^2=x^3
\]

has a cusp at the origin.

%%%%%%%%%%%%%%%%%%%%%%%%%%%%%%%%%%%%%%%%%%%%%%%%%%%%%%%%%%%%
\subsection{Node and Cusp}
\label{subsec:node-and-cusp}
%%%%%%%%%%%%%%%%%%%%%%%%%%%%%%%%%%%%%%%%%%%%%%%%%%%%%%%%%%%%

Two classical plane-curve singularities are:

\[
\boxed{
y^2=x^2(x+1)
}
\]

which has a node at the origin, and

\[
\boxed{
y^2=x^3
}
\]

which has a cusp.

A node locally resembles two branches crossing.

A cusp locally resembles a branch that comes to a sharp algebraic point.

Singularity theory studies such phenomena systematically.

%%%%%%%%%%%%%%%%%%%%%%%%%%%%%%%%%%%%%%%%%%%%%%%%%%%%%%%%%%%%
\subsection{Tangent Lines to Algebraic Curves}
\label{subsec:tangent-lines-algebraic-curves}
%%%%%%%%%%%%%%%%%%%%%%%%%%%%%%%%%%%%%%%%%%%%%%%%%%%%%%%%%%%%

Suppose

\[
f(x,y)=0
\]

and

\[
p=(a,b)
\]

is a nonsingular point.

The tangent line is determined by

\[
\boxed{
f_x(a,b)(x-a)
+
f_y(a,b)(y-b)
=
0.
}
\]

This follows from the first-order approximation of \(f\).

%%%%%%%%%%%%%%%%%%%%%%%%%%%%%%%%%%%%%%%%%%%%%%%%%%%%%%%%%%%%
\subsection{Worked Example: Tangent to a Circle}
\label{subsec:worked-algebraic-tangent-circle}
%%%%%%%%%%%%%%%%%%%%%%%%%%%%%%%%%%%%%%%%%%%%%%%%%%%%%%%%%%%%

\begin{workedexamplebox}{Tangent Line to an Algebraic Curve}

Consider

\[
f(x,y)=x^2+y^2-1.
\]

At

\[
p=(1,0),
\]

we have

\[
f_x=2x,
\qquad
f_y=2y.
\]

Thus

\[
f_x(1,0)=2,
\qquad
f_y(1,0)=0.
\]

The tangent equation is

\[
2(x-1)+0(y-0)=0.
\]

Therefore

\[
x=1.
\]

\begin{answerbox}
\[
\boxed{x=1}
\]
\end{answerbox}
\end{workedexamplebox}

%%%%%%%%%%%%%%%%%%%%%%%%%%%%%%%%%%%%%%%%%%%%%%%%%%%%%%%%%%%%
\subsection{Zariski Tangent Space}
\label{subsec:zariski-tangent-space}
%%%%%%%%%%%%%%%%%%%%%%%%%%%%%%%%%%%%%%%%%%%%%%%%%%%%%%%%%%%%

The tangent-space concept extends naturally to algebraic varieties.

Suppose

\[
X=V(f_1,\ldots,f_m)\subseteq\A^n_k
\]

and

\[
p\in X.
\]

The Zariski tangent space is described by the linear equations

\[
\boxed{
\sum_{j=1}^{n}
\frac{\partial f_i}{\partial x_j}(p)v_j=0,
\qquad
i=1,\ldots,m.
}
\]

Equivalently,

\[
T_pX
\]

is the kernel of the Jacobian matrix evaluated at \(p\).

%%%%%%%%%%%%%%%%%%%%%%%%%%%%%%%%%%%%%%%%%%%%%%%%%%%%%%%%%%%%
\subsection{The Jacobian Matrix}
\label{subsec:algebraic-jacobian-matrix}
%%%%%%%%%%%%%%%%%%%%%%%%%%%%%%%%%%%%%%%%%%%%%%%%%%%%%%%%%%%%

For

\[
f_1,\ldots,f_m,
\]

the Jacobian matrix is

\[
\boxed{
J=
\begin{pmatrix}
\dfrac{\partial f_1}{\partial x_1}
&
\cdots
&
\dfrac{\partial f_1}{\partial x_n}
\\[2mm]
\vdots
&
\ddots
&
\vdots
\\[2mm]
\dfrac{\partial f_m}{\partial x_1}
&
\cdots
&
\dfrac{\partial f_m}{\partial x_n}
\end{pmatrix}.
}
\]

At a point \(p\),

\[
\boxed{
T_pX=\ker J(p)
}
\]

for the affine algebraic set described by these equations.

The Jacobian criterion gives important methods for detecting smoothness.

%%%%%%%%%%%%%%%%%%%%%%%%%%%%%%%%%%%%%%%%%%%%%%%%%%%%%%%%%%%%
\subsection{Dimension and Singularities}
\label{subsec:dimension-singularities}
%%%%%%%%%%%%%%%%%%%%%%%%%%%%%%%%%%%%%%%%%%%%%%%%%%%%%%%%%%%%

At a smooth point of an irreducible variety \(X\),

\[
\boxed{
\dim T_pX=\dim X.
}
\]

At a singular point, the tangent space is typically too large:

\[
\boxed{
\dim T_pX>\dim X.
}
\]

For the cusp

\[
y^2=x^3,
\]

both first derivatives vanish at the origin.

Hence the linearized equations impose no restriction, so

\[
T_{(0,0)}X=\A^2.
\]

But the curve itself has dimension \(1\).

This excess tangent dimension detects the singularity.

%%%%%%%%%%%%%%%%%%%%%%%%%%%%%%%%%%%%%%%%%%%%%%%%%%%%%%%%%%%%
\subsection{Projective Algebraic Geometry}
\label{subsec:projective-algebraic-geometry}
%%%%%%%%%%%%%%%%%%%%%%%%%%%%%%%%%%%%%%%%%%%%%%%%%%%%%%%%%%%%

Affine space does not naturally contain points at infinity.

Projective space resolves this problem.

A point of projective \(n\)-space is written

\[
\boxed{
[X_0:X_1:\cdots:X_n].
}
\]

Two nonzero vectors represent the same point when

\[
\boxed{
[X_0:\cdots:X_n]
=
[\lambda X_0:\cdots:\lambda X_n]
}
\]

for every nonzero scalar

\[
\lambda\in k.
\]

Projective geometry was developed in
\cref{sec:projective-geometry}.

%%%%%%%%%%%%%%%%%%%%%%%%%%%%%%%%%%%%%%%%%%%%%%%%%%%%%%%%%%%%
\subsection{Homogeneous Polynomials}
\label{subsec:algebraic-homogeneous-polynomials}
%%%%%%%%%%%%%%%%%%%%%%%%%%%%%%%%%%%%%%%%%%%%%%%%%%%%%%%%%%%%

A polynomial

\[
F(X_0,\ldots,X_n)
\]

is homogeneous of degree \(d\) if every monomial has total degree \(d\).

For example,

\[
\boxed{
X^2+Y^2-Z^2
}
\]

is homogeneous of degree \(2\).

If

\[
F(P)=0,
\]

then

\[
F(\lambda P)
=
\lambda^dF(P)
=
0.
\]

Therefore the zero condition is well-defined on projective points.

%%%%%%%%%%%%%%%%%%%%%%%%%%%%%%%%%%%%%%%%%%%%%%%%%%%%%%%%%%%%
\subsection{Projective Algebraic Sets}
\label{subsec:projective-algebraic-sets}
%%%%%%%%%%%%%%%%%%%%%%%%%%%%%%%%%%%%%%%%%%%%%%%%%%%%%%%%%%%%

Projective algebraic sets are common zero sets of homogeneous polynomials in

\[
\mathbb{P}^n_k.
\]

For example,

\[
\boxed{
X^2+Y^2-Z^2=0
}
\]

defines a projective conic.

Setting

\[
Z=1
\]

recovers the affine equation

\[
x^2+y^2=1.
\]

The points satisfying

\[
Z=0
\]

are the points at infinity.

%%%%%%%%%%%%%%%%%%%%%%%%%%%%%%%%%%%%%%%%%%%%%%%%%%%%%%%%%%%%
\subsection{Projective Closure}
\label{subsec:projective-closure}
%%%%%%%%%%%%%%%%%%%%%%%%%%%%%%%%%%%%%%%%%%%%%%%%%%%%%%%%%%%%

An affine algebraic set can often be enlarged into a projective algebraic
set by homogenizing its equations.

For example,

\[
y=x^2
\]

can be written

\[
x^2-y=0.
\]

Homogenizing gives

\[
\boxed{
X^2-YZ=0.
}
\]

The projective curve

\[
V(X^2-YZ)
\subseteq\mathbb{P}^2
\]

contains the affine parabola together with its point at infinity.

%%%%%%%%%%%%%%%%%%%%%%%%%%%%%%%%%%%%%%%%%%%%%%%%%%%%%%%%%%%%
\subsection{Worked Example: Point at Infinity of a Parabola}
\label{subsec:worked-parabola-infinity}
%%%%%%%%%%%%%%%%%%%%%%%%%%%%%%%%%%%%%%%%%%%%%%%%%%%%%%%%%%%%

\begin{workedexamplebox}{Projective Closure of a Parabola}

Start with

\[
y=x^2.
\]

Its homogeneous equation is

\[
X^2=YZ.
\]

Points at infinity satisfy

\[
Z=0.
\]

Thus

\[
X^2=0,
\]

so

\[
X=0.
\]

The remaining coordinate \(Y\) must be nonzero, giving

\[
[0:1:0].
\]

\begin{answerbox}
\[
\boxed{
[0:1:0]
}
\]
\end{answerbox}
\end{workedexamplebox}

%%%%%%%%%%%%%%%%%%%%%%%%%%%%%%%%%%%%%%%%%%%%%%%%%%%%%%%%%%%%
\subsection{Why Projective Space Is Natural}
\label{subsec:why-projective-natural}
%%%%%%%%%%%%%%%%%%%%%%%%%%%%%%%%%%%%%%%%%%%%%%%%%%%%%%%%%%%%

Projective varieties often have better global properties than affine
varieties.

Projective space:

\begin{itemize}
    \item includes points at infinity;
    \item treats intersections more uniformly;
    \item provides the algebraic analogue of compactification;
    \item supports powerful intersection theorems;
    \item is central to modern algebraic geometry.
\end{itemize}

Many statements that require exceptions in affine space become cleaner in
projective space.

%%%%%%%%%%%%%%%%%%%%%%%%%%%%%%%%%%%%%%%%%%%%%%%%%%%%%%%%%%%%
\subsection{B\'ezout's Theorem}
\label{subsec:bezout-theorem}
%%%%%%%%%%%%%%%%%%%%%%%%%%%%%%%%%%%%%%%%%%%%%%%%%%%%%%%%%%%%

One of the classical motivations for projective geometry is intersection
theory.

\begin{theorem}[B\'ezout's Theorem --- Plane-Curve Form]
Over an algebraically closed field, two projective plane curves of degrees
\(m\) and \(n\) having no common irreducible component intersect in

\[
\boxed{mn}
\]

points when intersections are counted with appropriate multiplicities.
\end{theorem}

For example, two distinct projective lines intersect once:

\[
1\cdot1=1.
\]

A line and a conic intersect twice counting multiplicity:

\[
1\cdot2=2.
\]

%%%%%%%%%%%%%%%%%%%%%%%%%%%%%%%%%%%%%%%%%%%%%%%%%%%%%%%%%%%%
\subsection{Intersection Multiplicity}
\label{subsec:intersection-multiplicity}
%%%%%%%%%%%%%%%%%%%%%%%%%%%%%%%%%%%%%%%%%%%%%%%%%%%%%%%%%%%%

Two curves may meet tangentially rather than crossing transversely.

In such cases, the intersection point must be counted more than once.

For example, the line

\[
y=0
\]

is tangent to the parabola

\[
y=x^2
\]

at the origin.

Substitution gives

\[
x^2=0,
\]

which has a double root.

Thus the origin has intersection multiplicity \(2\).

Multiplicity allows B\'ezout's theorem to remain valid even for tangent
intersections.

%%%%%%%%%%%%%%%%%%%%%%%%%%%%%%%%%%%%%%%%%%%%%%%%%%%%%%%%%%%%
\subsection{Rational Functions}
\label{subsec:algebraic-rational-functions}
%%%%%%%%%%%%%%%%%%%%%%%%%%%%%%%%%%%%%%%%%%%%%%%%%%%%%%%%%%%%

On an irreducible affine variety \(X\), the coordinate ring

\[
k[X]
\]

is an integral domain.

Its field of fractions is called the \emph{function field} of \(X\):

\[
\boxed{
k(X)=\operatorname{Frac}(k[X]).
}
\]

Elements of

\[
k(X)
\]

are rational functions on \(X\).

For the affine line,

\[
\boxed{
k(\A^1)=k(x).
}
\]

%%%%%%%%%%%%%%%%%%%%%%%%%%%%%%%%%%%%%%%%%%%%%%%%%%%%%%%%%%%%
\subsection{Rational Maps}
\label{subsec:rational-maps}
%%%%%%%%%%%%%%%%%%%%%%%%%%%%%%%%%%%%%%%%%%%%%%%%%%%%%%%%%%%%

A rational map is allowed to be undefined on a proper algebraic subset.

Schematically,

\[
\boxed{
F:X\dashrightarrow Y.
}
\]

The dashed arrow indicates that the map need not be defined everywhere.

For example,

\[
\boxed{
F(x,y)=\left(\frac{x}{y},y\right)
}
\]

is undefined where

\[
y=0.
\]

Rational maps are essential because they capture geometry visible on dense
open subsets.

%%%%%%%%%%%%%%%%%%%%%%%%%%%%%%%%%%%%%%%%%%%%%%%%%%%%%%%%%%%%
\subsection{Birational Equivalence}
\label{subsec:birational-equivalence}
%%%%%%%%%%%%%%%%%%%%%%%%%%%%%%%%%%%%%%%%%%%%%%%%%%%%%%%%%%%%

\begin{definitionbox}{Birational Equivalence}
Two irreducible varieties \(X\) and \(Y\) are \emph{birationally equivalent}
if there exist rational maps

\[
X\dashrightarrow Y
\]

and

\[
Y\dashrightarrow X
\]

that are inverse to each other wherever both compositions are defined.
\end{definitionbox}

Birationally equivalent varieties have isomorphic function fields:

\[
\boxed{
k(X)\cong k(Y).
}
\]

Birational geometry studies varieties up to this weaker notion of
equivalence.

%%%%%%%%%%%%%%%%%%%%%%%%%%%%%%%%%%%%%%%%%%%%%%%%%%%%%%%%%%%%
\subsection{Rational Parametrization of a Circle}
\label{subsec:rational-parametrization-circle}
%%%%%%%%%%%%%%%%%%%%%%%%%%%%%%%%%%%%%%%%%%%%%%%%%%%%%%%%%%%%

Consider

\[
x^2+y^2=1.
\]

Draw a line of slope \(t\) through the point

\[
(-1,0).
\]

The line is

\[
y=t(x+1).
\]

Solving for the second intersection with the circle gives

\[
\boxed{
x=\frac{1-t^2}{1+t^2},
}
\]

\[
\boxed{
y=\frac{2t}{1+t^2}.
}
\]

Thus the circle is rationally parametrized.

This shows that the conic is birational to the affine or projective line.

%%%%%%%%%%%%%%%%%%%%%%%%%%%%%%%%%%%%%%%%%%%%%%%%%%%%%%%%%%%%
\subsection{Degree of a Variety}
\label{subsec:degree-variety}
%%%%%%%%%%%%%%%%%%%%%%%%%%%%%%%%%%%%%%%%%%%%%%%%%%%%%%%%%%%%

The degree of an algebraic variety generalizes the degree of a polynomial
and measures how the variety intersects generic linear spaces of
complementary dimension.

For a plane curve

\[
F(X,Y,Z)=0,
\]

the degree of the curve is the degree of the homogeneous polynomial \(F\)
when the equation is irreducible.

Examples include:

\[
\boxed{
\begin{array}{c|c}
\textbf{Curve} & \textbf{Degree}\\
\hline
\text{line} & 1\\
\text{conic} & 2\\
\text{cubic} & 3\\
\text{quartic} & 4
\end{array}
}
\]

%%%%%%%%%%%%%%%%%%%%%%%%%%%%%%%%%%%%%%%%%%%%%%%%%%%%%%%%%%%%
\subsection{Conics}
\label{subsec:algebraic-conics}
%%%%%%%%%%%%%%%%%%%%%%%%%%%%%%%%%%%%%%%%%%%%%%%%%%%%%%%%%%%%

A projective plane conic is defined by a homogeneous quadratic polynomial

\[
\boxed{
F(X,Y,Z)=0.
}
\]

Over an algebraically closed field, every nonsingular projective conic is
isomorphic to

\[
\mathbb{P}^1.
\]

Over non-algebraically-closed fields, arithmetic questions become important.

For example, a conic over \(\Q\) may or may not possess a rational point.

This connects algebraic geometry with Diophantine Geometry.

%%%%%%%%%%%%%%%%%%%%%%%%%%%%%%%%%%%%%%%%%%%%%%%%%%%%%%%%%%%%
\subsection{Cubic Curves and Elliptic Curves}
\label{subsec:algebraic-cubic-curves}
%%%%%%%%%%%%%%%%%%%%%%%%%%%%%%%%%%%%%%%%%%%%%%%%%%%%%%%%%%%%

A smooth projective cubic curve has genus \(1\).

When supplied with a distinguished rational point, it becomes an elliptic
curve.

A common affine equation is

\[
\boxed{
y^2=x^3+Ax+B
}
\]

with

\[
4A^3+27B^2\neq0
\]

in characteristic not equal to \(2\) or \(3\).

Elliptic curves were introduced arithmetically in
\cref{sec:diophantine-geometry} and
\cref{sec:arithmetic-geometry}.

Here the emphasis is their geometric structure.

%%%%%%%%%%%%%%%%%%%%%%%%%%%%%%%%%%%%%%%%%%%%%%%%%%%%%%%%%%%%
\subsection{Genus}
\label{subsec:algebraic-genus}
%%%%%%%%%%%%%%%%%%%%%%%%%%%%%%%%%%%%%%%%%%%%%%%%%%%%%%%%%%%%

The genus is a fundamental invariant of algebraic curves.

For a smooth projective complex curve, it agrees with the topological genus
of the corresponding compact Riemann surface.

Very roughly,

\[
\boxed{
\begin{array}{c|c}
g=0 & \text{rational-curve behavior}\\
g=1 & \text{elliptic-curve behavior}\\
g\geq2 & \text{higher-genus behavior}
\end{array}
}
\]

For a smooth projective plane curve of degree \(d\),

\[
\boxed{
g=
\frac{(d-1)(d-2)}{2}.
}
\]

Thus a smooth conic has genus \(0\), while a smooth cubic has genus \(1\).

%%%%%%%%%%%%%%%%%%%%%%%%%%%%%%%%%%%%%%%%%%%%%%%%%%%%%%%%%%%%
\subsection{Worked Example: Genus of a Plane Quartic}
\label{subsec:worked-genus-quartic}
%%%%%%%%%%%%%%%%%%%%%%%%%%%%%%%%%%%%%%%%%%%%%%%%%%%%%%%%%%%%

\begin{workedexamplebox}{Genus of a Smooth Plane Quartic}

For a smooth projective plane curve of degree

\[
d=4,
\]

the genus is

\[
g
=
\frac{(d-1)(d-2)}{2}.
\]

Therefore

\[
g
=
\frac{(3)(2)}{2}
=
3.
\]

\begin{answerbox}
\[
\boxed{g=3}
\]
\end{answerbox}
\end{workedexamplebox}

%%%%%%%%%%%%%%%%%%%%%%%%%%%%%%%%%%%%%%%%%%%%%%%%%%%%%%%%%%%%
\subsection{Divisors}
\label{subsec:algebraic-divisors}
%%%%%%%%%%%%%%%%%%%%%%%%%%%%%%%%%%%%%%%%%%%%%%%%%%%%%%%%%%%%

On a smooth algebraic curve, a divisor is a formal finite sum of points:

\[
\boxed{
D=\sum_P n_P[P],
}
\]

where

\[
n_P\in\Z
\]

and all but finitely many coefficients are zero.

The degree is

\[
\boxed{
\deg D=\sum_P n_P.
}
\]

Divisors encode zeros and poles of rational functions and play a central
role in the geometry of curves.

%%%%%%%%%%%%%%%%%%%%%%%%%%%%%%%%%%%%%%%%%%%%%%%%%%%%%%%%%%%%
\subsection{Principal Divisors}
\label{subsec:principal-divisors}
%%%%%%%%%%%%%%%%%%%%%%%%%%%%%%%%%%%%%%%%%%%%%%%%%%%%%%%%%%%%

A nonzero rational function

\[
f\in k(X)
\]

has zeros and poles.

Its associated divisor is

\[
\boxed{
\operatorname{div}(f)
=
\sum_P
\operatorname{ord}_P(f)[P].
}
\]

Here

\[
\operatorname{ord}_P(f)
\]

records the order of vanishing or pole of \(f\) at \(P\).

Such divisors are called \emph{principal divisors}.

%%%%%%%%%%%%%%%%%%%%%%%%%%%%%%%%%%%%%%%%%%%%%%%%%%%%%%%%%%%%
\subsection{Divisor Classes and the Picard Group}
\label{subsec:picard-group}
%%%%%%%%%%%%%%%%%%%%%%%%%%%%%%%%%%%%%%%%%%%%%%%%%%%%%%%%%%%%

Two divisors are linearly equivalent if their difference is principal.

The resulting divisor classes form a group closely related to the
\emph{Picard group}

\[
\boxed{
\operatorname{Pic}(X).
}
\]

In modern language, the Picard group may also be described using isomorphism
classes of line bundles.

This is one of the first major bridges from elementary algebraic varieties
to sheaf-theoretic geometry.

%%%%%%%%%%%%%%%%%%%%%%%%%%%%%%%%%%%%%%%%%%%%%%%%%%%%%%%%%%%%
\subsection{Riemann--Roch Philosophy}
\label{subsec:riemann-roch-philosophy}
%%%%%%%%%%%%%%%%%%%%%%%%%%%%%%%%%%%%%%%%%%%%%%%%%%%%%%%%%%%%

Given a divisor \(D\) on a smooth projective curve, one may ask how many
rational functions are allowed to have poles bounded by \(D\).

The Riemann--Roch theorem relates this analytic-algebraic quantity to the
degree and genus of the curve.

\begin{theorem}[Riemann--Roch Theorem for Curves]
For a divisor \(D\) on a smooth projective curve of genus \(g\),

\[
\boxed{
\ell(D)-\ell(K-D)
=
\deg D+1-g,
}
\]

where \(K\) is a canonical divisor.
\end{theorem}

This theorem is one of the foundational results of algebraic geometry.

%%%%%%%%%%%%%%%%%%%%%%%%%%%%%%%%%%%%%%%%%%%%%%%%%%%%%%%%%%%%
\subsection{Why Local Information Matters}
\label{subsec:local-information-algebraic-geometry}
%%%%%%%%%%%%%%%%%%%%%%%%%%%%%%%%%%%%%%%%%%%%%%%%%%%%%%%%%%%%

A variety may look very different near different points.

To understand the geometry near a point \(p\), algebraic geometers study
functions defined near \(p\).

This leads to the \emph{local ring}

\[
\boxed{
\mathcal{O}_{X,p}.
}
\]

The symbol

\[
\mathcal{O}_{X,p}
\]

is read ``O sub \(X,p\).''

It records regular functions defined in some neighborhood of \(p\), with
functions identified when they agree near \(p\).

%%%%%%%%%%%%%%%%%%%%%%%%%%%%%%%%%%%%%%%%%%%%%%%%%%%%%%%%%%%%
\subsection{Local Rings}
\label{subsec:geometry-local-rings}
%%%%%%%%%%%%%%%%%%%%%%%%%%%%%%%%%%%%%%%%%%%%%%%%%%%%%%%%%%%%

A local ring has a unique maximal ideal.

For a variety \(X\) and point \(p\), the maximal ideal of

\[
\mathcal{O}_{X,p}
\]

consists of germs of regular functions vanishing at \(p\).

The local ring contains detailed information about:

\begin{itemize}
    \item smoothness;
    \item dimension;
    \item multiplicity;
    \item tangent directions;
    \item singularities.
\end{itemize}

Thus local algebra becomes a microscope for geometry.

%%%%%%%%%%%%%%%%%%%%%%%%%%%%%%%%%%%%%%%%%%%%%%%%%%%%%%%%%%%%
\subsection{From Varieties to Schemes}
\label{subsec:varieties-to-schemes}
%%%%%%%%%%%%%%%%%%%%%%%%%%%%%%%%%%%%%%%%%%%%%%%%%%%%%%%%%%%%

Classical varieties focus primarily on polynomial solution sets over fields.

Modern algebraic geometry broadens the viewpoint.

Instead of beginning with solutions to equations, one begins with a
commutative ring.

\begin{definitionbox}{Spectrum of a Ring}
For a commutative ring \(R\), the \emph{spectrum} of \(R\) is

\[
\boxed{
\Spec(R)
=
\{
\mathfrak p:
\mathfrak p
\text{ is a prime ideal of }R
\}.
}
\]
\end{definitionbox}

The notation

\[
\Spec(R)
\]

is read ``spectrum of \(R\)'' or ``Spec \(R\).''

%%%%%%%%%%%%%%%%%%%%%%%%%%%%%%%%%%%%%%%%%%%%%%%%%%%%%%%%%%%%
\subsection{Prime Ideals as Points}
\label{subsec:prime-ideals-as-points}
%%%%%%%%%%%%%%%%%%%%%%%%%%%%%%%%%%%%%%%%%%%%%%%%%%%%%%%%%%%%

In classical affine geometry, ordinary points correspond to maximal ideals
over an algebraically closed field.

Scheme theory enlarges the space by allowing every prime ideal to act as a
point.

Thus

\[
\boxed{
\text{points of }\Spec(R)
=
\text{prime ideals of }R.
}
\]

Maximal ideals are still present, but nonmaximal prime ideals provide
additional generic information.

%%%%%%%%%%%%%%%%%%%%%%%%%%%%%%%%%%%%%%%%%%%%%%%%%%%%%%%%%%%%
\subsection{Example: Spectrum of the Integers}
\label{subsec:spectrum-integers}
%%%%%%%%%%%%%%%%%%%%%%%%%%%%%%%%%%%%%%%%%%%%%%%%%%%%%%%%%%%%

The prime ideals of \(\Z\) are

\[
(0)
\]

and

\[
(p)
\]

for prime numbers \(p\).

Therefore

\[
\boxed{
\Spec(\Z)
=
\{
(0),(2),(3),(5),(7),\ldots
\}.
}
\]

This gives a geometric interpretation of prime numbers as points of an
arithmetic space.

The point

\[
(0)
\]

acts as a generic point.

This viewpoint is fundamental in Arithmetic Geometry.

%%%%%%%%%%%%%%%%%%%%%%%%%%%%%%%%%%%%%%%%%%%%%%%%%%%%%%%%%%%%
\subsection{Zariski Topology on the Spectrum}
\label{subsec:zariski-spectrum}
%%%%%%%%%%%%%%%%%%%%%%%%%%%%%%%%%%%%%%%%%%%%%%%%%%%%%%%%%%%%

For an ideal

\[
I\subseteq R,
\]

define

\[
\boxed{
V(I)
=
\{
\mathfrak p\in\Spec(R):
I\subseteq\mathfrak p
\}.
}
\]

These sets are declared closed.

Thus

\[
\Spec(R)
\]

becomes a topological space.

The basic open sets are

\[
\boxed{
D(f)
=
\{
\mathfrak p:
f\notin\mathfrak p
\}.
}
\]

This generalizes the Zariski topology of affine varieties.

%%%%%%%%%%%%%%%%%%%%%%%%%%%%%%%%%%%%%%%%%%%%%%%%%%%%%%%%%%%%
\subsection{Affine Schemes}
\label{subsec:affine-schemes}
%%%%%%%%%%%%%%%%%%%%%%%%%%%%%%%%%%%%%%%%%%%%%%%%%%%%%%%%%%%%

A topological space alone does not retain enough algebraic information.

The spectrum is therefore equipped with a sheaf of rings.

\begin{definitionbox}{Affine Scheme}
An \emph{affine scheme} is a locally ringed space of the form

\[
\boxed{
\bigl(
\Spec(R),
\mathcal{O}_{\Spec(R)}
\bigr)
}
\]

for a commutative ring \(R\).
\end{definitionbox}

Usually this is abbreviated simply as

\[
\boxed{
\Spec(R).
}
\]

%%%%%%%%%%%%%%%%%%%%%%%%%%%%%%%%%%%%%%%%%%%%%%%%%%%%%%%%%%%%
\subsection{What Is a Sheaf?}
\label{subsec:algebraic-geometry-sheaves}
%%%%%%%%%%%%%%%%%%%%%%%%%%%%%%%%%%%%%%%%%%%%%%%%%%%%%%%%%%%%

A sheaf is a mathematical device for consistently assigning local data to
open subsets of a space.

For each open set \(U\), a sheaf assigns some collection

\[
\mathcal{F}(U).
\]

For a sheaf of functions, these might be functions defined on \(U\).

The sheaf includes rules for:

\begin{itemize}
    \item restricting data to smaller open sets;
    \item checking when local data agree;
    \item gluing compatible local data into global data.
\end{itemize}

Schematically,

\[
\boxed{
\text{local data}
+
\text{compatibility}
\longrightarrow
\text{global data}.
}
\]

%%%%%%%%%%%%%%%%%%%%%%%%%%%%%%%%%%%%%%%%%%%%%%%%%%%%%%%%%%%%
\subsection{The Structure Sheaf}
\label{subsec:structure-sheaf}
%%%%%%%%%%%%%%%%%%%%%%%%%%%%%%%%%%%%%%%%%%%%%%%%%%%%%%%%%%%%

The sheaf attached to a scheme is called the \emph{structure sheaf}.

It is commonly denoted

\[
\boxed{
\mathcal{O}_X.
}
\]

For an affine scheme

\[
X=\Spec(R),
\]

the global sections satisfy

\[
\boxed{
\Gamma(X,\mathcal{O}_X)\cong R.
}
\]

Thus the original ring can be recovered from its affine scheme.

%%%%%%%%%%%%%%%%%%%%%%%%%%%%%%%%%%%%%%%%%%%%%%%%%%%%%%%%%%%%
\subsection{Why Schemes Are More Powerful}
\label{subsec:why-schemes}
%%%%%%%%%%%%%%%%%%%%%%%%%%%%%%%%%%%%%%%%%%%%%%%%%%%%%%%%%%%%

Schemes allow algebraic geometry to study phenomena invisible to ordinary
point sets.

They can encode:

\begin{itemize}
    \item nilpotent elements;
    \item multiplicities;
    \item infinitesimal structure;
    \item geometry over arbitrary commutative rings;
    \item families of varieties;
    \item arithmetic information;
    \item reduction modulo primes.
\end{itemize}

Recall

\[
V(x)=V(x^2).
\]

Classical zero sets cannot distinguish these equations.

Scheme structures can.

%%%%%%%%%%%%%%%%%%%%%%%%%%%%%%%%%%%%%%%%%%%%%%%%%%%%%%%%%%%%
\subsection{Nilpotents and Infinitesimal Thickness}
\label{subsec:nilpotents-infinitesimal-thickness}
%%%%%%%%%%%%%%%%%%%%%%%%%%%%%%%%%%%%%%%%%%%%%%%%%%%%%%%%%%%%

Compare

\[
k[x]/(x)
\]

with

\[
k[x]/(x^2).
\]

Both have the same ordinary geometric point

\[
x=0.
\]

But in

\[
k[x]/(x^2),
\]

the class of \(x\) is nonzero while

\[
x^2=0.
\]

This nilpotent element records an infinitesimal thickening of the point.

Schematically,

\[
\boxed{
(x)
\quad\text{ordinary point},
}
\]

while

\[
\boxed{
(x^2)
\quad\text{double or infinitesimally thickened point}.
}
\]

%%%%%%%%%%%%%%%%%%%%%%%%%%%%%%%%%%%%%%%%%%%%%%%%%%%%%%%%%%%%
\subsection{Schemes as Locally Affine Spaces}
\label{subsec:schemes-locally-affine}
%%%%%%%%%%%%%%%%%%%%%%%%%%%%%%%%%%%%%%%%%%%%%%%%%%%%%%%%%%%%

A general scheme is built by gluing affine schemes together.

This parallels the construction of manifolds from coordinate charts.

\[
\boxed{
\begin{array}{c|c}
\textbf{Differential Geometry}
&
\textbf{Algebraic Geometry}
\\
\hline
\text{Euclidean chart}
&
\text{affine scheme}
\\
\text{smooth functions}
&
\text{regular functions}
\\
\text{smooth manifold}
&
\text{scheme}
\\
\text{local coordinates}
&
\text{local rings}
\end{array}
}
\]

The analogy is useful, although the categories have important differences.

%%%%%%%%%%%%%%%%%%%%%%%%%%%%%%%%%%%%%%%%%%%%%%%%%%%%%%%%%%%%
\subsection{Generic Points}
\label{subsec:generic-points}
%%%%%%%%%%%%%%%%%%%%%%%%%%%%%%%%%%%%%%%%%%%%%%%%%%%%%%%%%%%%

One unusual feature of schemes is the existence of generic points.

For example, in

\[
\Spec(k[x]),
\]

the prime ideal

\[
(0)
\]

corresponds to the generic point of the entire affine line.

Its closure is the whole space.

Thus a single scheme-theoretic point can represent an entire irreducible
subvariety.

%%%%%%%%%%%%%%%%%%%%%%%%%%%%%%%%%%%%%%%%%%%%%%%%%%%%%%%%%%%%
\subsection{Residue Fields}
\label{subsec:residue-fields}
%%%%%%%%%%%%%%%%%%%%%%%%%%%%%%%%%%%%%%%%%%%%%%%%%%%%%%%%%%%%

Every point

\[
\mathfrak p\in\Spec(R)
\]

has an associated field called its \emph{residue field}.

It is denoted

\[
\boxed{
\kappa(\mathfrak p).
}
\]

The Greek letter

\[
\kappa
\]

is read ``kappa.''

One construction is

\[
\boxed{
\kappa(\mathfrak p)
=
\operatorname{Frac}(R/\mathfrak p)
}
\]

after localizing appropriately; equivalently,

\[
\kappa(\mathfrak p)
=
R_{\mathfrak p}/\mathfrak pR_{\mathfrak p}.
\]

Residue fields record the field naturally attached to a scheme point.

%%%%%%%%%%%%%%%%%%%%%%%%%%%%%%%%%%%%%%%%%%%%%%%%%%%%%%%%%%%%
\subsection{Morphisms of Schemes}
\label{subsec:morphisms-of-schemes}
%%%%%%%%%%%%%%%%%%%%%%%%%%%%%%%%%%%%%%%%%%%%%%%%%%%%%%%%%%%%

A ring homomorphism

\[
\varphi:R\to S
\]

induces a morphism in the opposite direction:

\[
\boxed{
\Spec(S)\longrightarrow\Spec(R).
}
\]

A prime ideal

\[
\mathfrak q\subseteq S
\]

is sent to its inverse image

\[
\boxed{
\varphi^{-1}(\mathfrak q)\subseteq R.
}
\]

Again,

\[
\boxed{
\text{algebra}
\quad\text{and}\quad
\text{geometry}
}
\]

move contravariantly.

%%%%%%%%%%%%%%%%%%%%%%%%%%%%%%%%%%%%%%%%%%%%%%%%%%%%%%%%%%%%
\subsection{Fibers of Morphisms}
\label{subsec:scheme-fibers}
%%%%%%%%%%%%%%%%%%%%%%%%%%%%%%%%%%%%%%%%%%%%%%%%%%%%%%%%%%%%

Suppose

\[
f:X\to Y
\]

is a morphism of schemes and

\[
y\in Y.
\]

The fiber over \(y\) is, schematically,

\[
\boxed{
X_y
=
X\times_Y\Spec(\kappa(y)).
}
\]

Fibers allow a family of geometric objects to vary over a base space.

This becomes especially powerful when

\[
Y=\Spec(\Z).
\]

Then different fibers correspond to different prime characteristics.

%%%%%%%%%%%%%%%%%%%%%%%%%%%%%%%%%%%%%%%%%%%%%%%%%%%%%%%%%%%%
\subsection{Arithmetic Families}
\label{subsec:arithmetic-families}
%%%%%%%%%%%%%%%%%%%%%%%%%%%%%%%%%%%%%%%%%%%%%%%%%%%%%%%%%%%%

Consider an equation with integer coefficients, such as

\[
y^2=x^3-x+1.
\]

It can be viewed not merely as one curve but as a family over

\[
\Spec(\Z).
\]

The generic fiber lives over

\[
\Q,
\]

while the fiber over a prime \(p\) is obtained by reducing the equation
modulo \(p\):

\[
\boxed{
y^2=x^3-x+1
\pmod p.
}
\]

This geometric viewpoint was introduced in
\cref{sec:arithmetic-geometry}.

%%%%%%%%%%%%%%%%%%%%%%%%%%%%%%%%%%%%%%%%%%%%%%%%%%%%%%%%%%%%
\subsection{Base Change}
\label{subsec:algebraic-base-change}
%%%%%%%%%%%%%%%%%%%%%%%%%%%%%%%%%%%%%%%%%%%%%%%%%%%%%%%%%%%%

Geometry can change when the underlying field or base ring changes.

For example, an algebraic variety over \(\R\) may be extended to \(\C\).

Schematically,

\[
\boxed{
X_{\C}
=
X_{\R}\times_{\Spec(\R)}\Spec(\C).
}
\]

This operation is called \emph{base change}.

Base change allows one geometric object to be studied over different fields
while retaining its algebraic origin.

%%%%%%%%%%%%%%%%%%%%%%%%%%%%%%%%%%%%%%%%%%%%%%%%%%%%%%%%%%%%
\subsection{Products and Fiber Products}
\label{subsec:algebraic-fiber-products}
%%%%%%%%%%%%%%%%%%%%%%%%%%%%%%%%%%%%%%%%%%%%%%%%%%%%%%%%%%%%

The categorical product in algebraic geometry is the \emph{fiber product}.

Given

\[
X\to S
\]

and

\[
Y\to S,
\]

their fiber product is written

\[
\boxed{
X\times_S Y.
}
\]

It represents pairs of points of \(X\) and \(Y\) lying over the same point
of \(S\), together with the appropriate scheme structure.

Fiber products are fundamental for:

\begin{itemize}
    \item base change;
    \item intersections;
    \item families;
    \item moduli problems.
\end{itemize}

%%%%%%%%%%%%%%%%%%%%%%%%%%%%%%%%%%%%%%%%%%%%%%%%%%%%%%%%%%%%
\subsection{Algebraic Geometry and Category Theory}
\label{subsec:algebraic-geometry-category-theory}
%%%%%%%%%%%%%%%%%%%%%%%%%%%%%%%%%%%%%%%%%%%%%%%%%%%%%%%%%%%%

Modern algebraic geometry is strongly categorical.

The construction

\[
R\longmapsto\Spec(R)
\]

reverses arrows.

Thus affine schemes are dual to commutative rings in a precise categorical
sense:

\[
\boxed{
\operatorname{Hom}_{\mathrm{Sch}}
\bigl(
\Spec(S),\Spec(R)
\bigr)
\cong
\operatorname{Hom}_{\mathrm{CRing}}(R,S).
}
\]

This explains why Category Theory is deeply embedded in modern algebraic
geometry.

%%%%%%%%%%%%%%%%%%%%%%%%%%%%%%%%%%%%%%%%%%%%%%%%%%%%%%%%%%%%
\subsection{Algebraic Geometry and Complex Geometry}
\label{subsec:algebraic-geometry-complex-geometry}
%%%%%%%%%%%%%%%%%%%%%%%%%%%%%%%%%%%%%%%%%%%%%%%%%%%%%%%%%%%%

When polynomial equations are studied over

\[
\C,
\]

their solution sets can often be viewed simultaneously as algebraic and
complex-analytic spaces.

For example, a smooth complex algebraic curve gives a Riemann surface.

This creates a powerful correspondence among:

\[
\boxed{
\text{algebraic curves over }\C,
}
\]

\[
\boxed{
\text{compact Riemann surfaces},
}
\]

and

\[
\boxed{
\text{one-dimensional complex manifolds}.
}
\]

This relationship is developed further in
\cref{sec:complex-geometry}.

%%%%%%%%%%%%%%%%%%%%%%%%%%%%%%%%%%%%%%%%%%%%%%%%%%%%%%%%%%%%
\subsection{Algebraic Geometry and Number Theory}
\label{subsec:algebraic-geometry-number-theory}
%%%%%%%%%%%%%%%%%%%%%%%%%%%%%%%%%%%%%%%%%%%%%%%%%%%%%%%%%%%%

Number theory asks for solutions over fields such as

\[
\Q
\]

and finite fields

\[
\F_p.
\]

Algebraic geometry supplies the geometric objects on which those arithmetic
questions are posed.

For example,

\[
E:y^2=x^3+Ax+B
\]

defines an algebraic curve.

Number theory then asks about

\[
E(\Q),
\qquad
E(\F_p),
\qquad
E(\Q_p).
\]

This interaction creates Diophantine Geometry and Arithmetic Geometry.

%%%%%%%%%%%%%%%%%%%%%%%%%%%%%%%%%%%%%%%%%%%%%%%%%%%%%%%%%%%%
\subsection{Algebraic Geometry and Topology}
\label{subsec:algebraic-geometry-topology}
%%%%%%%%%%%%%%%%%%%%%%%%%%%%%%%%%%%%%%%%%%%%%%%%%%%%%%%%%%%%

Complex algebraic varieties have underlying topological spaces.

Topological invariants can therefore reveal algebraic information.

Important connections include:

\begin{itemize}
    \item homology;
    \item cohomology;
    \item fundamental groups;
    \item characteristic classes;
    \item intersection theory.
\end{itemize}

Modern algebraic geometry also develops purely algebraic analogues of
topological invariants, such as \'etale cohomology.

%%%%%%%%%%%%%%%%%%%%%%%%%%%%%%%%%%%%%%%%%%%%%%%%%%%%%%%%%%%%
\subsection{Algebraic Geometry and Differential Geometry}
\label{subsec:algebraic-differential-geometry}
%%%%%%%%%%%%%%%%%%%%%%%%%%%%%%%%%%%%%%%%%%%%%%%%%%%%%%%%%%%%

A smooth complex algebraic variety may also carry differential-geometric
structures.

This permits tools such as:

\begin{itemize}
    \item differential forms;
    \item Hermitian metrics;
    \item curvature;
    \item Hodge theory;
    \item K\"ahler geometry.
\end{itemize}

These interactions become central in Complex Geometry.

%%%%%%%%%%%%%%%%%%%%%%%%%%%%%%%%%%%%%%%%%%%%%%%%%%%%%%%%%%%%
\subsection{Moduli Spaces}
\label{subsec:moduli-spaces-preview}
%%%%%%%%%%%%%%%%%%%%%%%%%%%%%%%%%%%%%%%%%%%%%%%%%%%%%%%%%%%%

A \emph{moduli problem} asks for a geometric space whose points represent
other geometric objects.

For example, one might seek a space whose points represent:

\begin{itemize}
    \item algebraic curves of a fixed genus;
    \item elliptic curves;
    \item vector bundles;
    \item line bundles;
    \item hypersurfaces.
\end{itemize}

Schematically,

\[
\boxed{
\text{point in moduli space}
\quad\longleftrightarrow\quad
\text{isomorphism class of geometric object}.
}
\]

Moduli theory is one of the major advanced areas of algebraic geometry.

%%%%%%%%%%%%%%%%%%%%%%%%%%%%%%%%%%%%%%%%%%%%%%%%%%%%%%%%%%%%
\subsection{The Moduli of Elliptic Curves}
\label{subsec:moduli-elliptic-curves}
%%%%%%%%%%%%%%%%%%%%%%%%%%%%%%%%%%%%%%%%%%%%%%%%%%%%%%%%%%%%

Over an algebraically closed field of suitable characteristic, elliptic
curves are classified up to isomorphism by an invariant called the
\(j\)-invariant.

It is written

\[
\boxed{
j(E).
}
\]

For a Weierstrass equation

\[
y^2=x^3+Ax+B,
\]

one standard formula is

\[
\boxed{
j
=
1728
\frac{4A^3}
{4A^3+27B^2}.
}
\]

Two elliptic curves over an algebraically closed field of characteristic not
\(2\) or \(3\) are isomorphic if and only if they have the same
\(j\)-invariant.

%%%%%%%%%%%%%%%%%%%%%%%%%%%%%%%%%%%%%%%%%%%%%%%%%%%%%%%%%%%%
\subsection{Line Bundles}
\label{subsec:algebraic-line-bundles}
%%%%%%%%%%%%%%%%%%%%%%%%%%%%%%%%%%%%%%%%%%%%%%%%%%%%%%%%%%%%

A line bundle assigns a one-dimensional vector space to every point of a
variety in a locally compatible way.

Line bundles are algebraic analogues of geometric families of lines.

They are important because they encode:

\begin{itemize}
    \item divisors;
    \item embeddings into projective space;
    \item sections;
    \item positivity;
    \item intersection theory.
\end{itemize}

The Picard group organizes line bundles up to isomorphism.

%%%%%%%%%%%%%%%%%%%%%%%%%%%%%%%%%%%%%%%%%%%%%%%%%%%%%%%%%%%%
\subsection{Projective Embeddings}
\label{subsec:projective-embeddings}
%%%%%%%%%%%%%%%%%%%%%%%%%%%%%%%%%%%%%%%%%%%%%%%%%%%%%%%%%%%%

Line bundles and their global sections can be used to map varieties into
projective space.

A sufficiently positive line bundle may produce an embedding

\[
\boxed{
X\hookrightarrow\mathbb{P}^N.
}
\]

This allows an abstract variety to be represented by homogeneous polynomial
equations.

Thus projective geometry remains central even in the abstract theory.

%%%%%%%%%%%%%%%%%%%%%%%%%%%%%%%%%%%%%%%%%%%%%%%%%%%%%%%%%%%%
\subsection{Algebraic Differential Forms}
\label{subsec:algebraic-differential-forms}
%%%%%%%%%%%%%%%%%%%%%%%%%%%%%%%%%%%%%%%%%%%%%%%%%%%%%%%%%%%%

Differential forms also have an algebraic counterpart.

For a \(k\)-algebra \(A\), one defines the module of K\"ahler differentials

\[
\boxed{
\Omega_{A/k}.
}
\]

The universal derivation is

\[
\boxed{
d:A\to\Omega_{A/k}.
}
\]

This algebraic construction captures infinitesimal behavior without relying
on analytic limits.

K\"ahler differentials are fundamental in the study of smoothness, tangent
spaces, and deformation theory.

%%%%%%%%%%%%%%%%%%%%%%%%%%%%%%%%%%%%%%%%%%%%%%%%%%%%%%%%%%%%
\subsection{Tangent Spaces from Local Algebra}
\label{subsec:tangent-space-local-algebra}
%%%%%%%%%%%%%%%%%%%%%%%%%%%%%%%%%%%%%%%%%%%%%%%%%%%%%%%%%%%%

Let \(p\) be a point with maximal ideal

\[
\mathfrak m_p.
\]

The cotangent space at \(p\) can be described algebraically by

\[
\boxed{
\mathfrak m_p/\mathfrak m_p^2.
}
\]

The tangent space is its dual:

\[
\boxed{
T_pX
\cong
\operatorname{Hom}_k
\left(
\mathfrak m_p/\mathfrak m_p^2,
k
\right)
}
\]

under suitable hypotheses.

This formula shows how tangent directions are encoded entirely by local
ring structure.

%%%%%%%%%%%%%%%%%%%%%%%%%%%%%%%%%%%%%%%%%%%%%%%%%%%%%%%%%%%%
\subsection{Smoothness as an Algebraic Property}
\label{subsec:algebraic-smoothness}
%%%%%%%%%%%%%%%%%%%%%%%%%%%%%%%%%%%%%%%%%%%%%%%%%%%%%%%%%%%%

Differential geometry defines smoothness using derivatives.

Algebraic geometry can detect analogous behavior using local rings,
Jacobians, and K\"ahler differentials.

For a variety \(X\), a point \(p\) is smooth when its local algebra has the
expected regularity.

For varieties over suitable fields, one important criterion is

\[
\boxed{
\dim T_pX=\dim_pX.
}
\]

When the tangent space is too large, the point is singular.

%%%%%%%%%%%%%%%%%%%%%%%%%%%%%%%%%%%%%%%%%%%%%%%%%%%%%%%%%%%%
\subsection{Resolution of Singularities}
\label{subsec:resolution-singularities}
%%%%%%%%%%%%%%%%%%%%%%%%%%%%%%%%%%%%%%%%%%%%%%%%%%%%%%%%%%%%

A singular variety can sometimes be replaced by a smooth variety that
retains much of its geometry.

Schematically,

\[
\boxed{
\widetilde{X}
\longrightarrow
X,
}
\]

where

\[
\widetilde{X}
\]

is smooth and the map is an isomorphism away from the singular region.

This process is called \emph{resolution of singularities}.

In characteristic zero, Hironaka's theorem guarantees resolution for
algebraic varieties.

%%%%%%%%%%%%%%%%%%%%%%%%%%%%%%%%%%%%%%%%%%%%%%%%%%%%%%%%%%%%
\subsection{Blowups}
\label{subsec:algebraic-blowups}
%%%%%%%%%%%%%%%%%%%%%%%%%%%%%%%%%%%%%%%%%%%%%%%%%%%%%%%%%%%%

One of the fundamental tools for resolving singularities is the
\emph{blowup}.

Very roughly, blowing up a point replaces the point by the collection of
directions passing through it.

For a point in a surface, the exceptional set often resembles

\[
\boxed{
\mathbb{P}^1.
}
\]

Blowups convert directional information hidden at a singular point into
visible geometric structure.

%%%%%%%%%%%%%%%%%%%%%%%%%%%%%%%%%%%%%%%%%%%%%%%%%%%%%%%%%%%%
\subsection{Birational Geometry}
\label{subsec:birational-geometry-preview}
%%%%%%%%%%%%%%%%%%%%%%%%%%%%%%%%%%%%%%%%%%%%%%%%%%%%%%%%%%%%

Birational geometry studies varieties while allowing modifications that do
not change their rational function fields.

Typical operations include:

\begin{itemize}
    \item blowups;
    \item blowdowns;
    \item resolutions of singularities;
    \item rational maps.
\end{itemize}

The broad classification problem asks:

\[
\boxed{
\text{How can algebraic varieties be classified up to birational
equivalence?}
}
\]

In higher dimensions this leads to the Minimal Model Program.

%%%%%%%%%%%%%%%%%%%%%%%%%%%%%%%%%%%%%%%%%%%%%%%%%%%%%%%%%%%%
\subsection{The Minimal Model Program}
\label{subsec:minimal-model-program-preview}
%%%%%%%%%%%%%%%%%%%%%%%%%%%%%%%%%%%%%%%%%%%%%%%%%%%%%%%%%%%%

The Minimal Model Program attempts to simplify algebraic varieties within
their birational equivalence classes.

Its philosophy is analogous to finding canonical or simplest
representatives.

Important concepts include:

\begin{itemize}
    \item canonical divisors;
    \item contractions;
    \item flips;
    \item minimal models;
    \item Fano varieties.
\end{itemize}

The full theory belongs to advanced birational algebraic geometry.

%%%%%%%%%%%%%%%%%%%%%%%%%%%%%%%%%%%%%%%%%%%%%%%%%%%%%%%%%%%%
\subsection{Cohomology in Algebraic Geometry}
\label{subsec:algebraic-cohomology-preview}
%%%%%%%%%%%%%%%%%%%%%%%%%%%%%%%%%%%%%%%%%%%%%%%%%%%%%%%%%%%%

Sheaf cohomology measures obstructions to gluing local information globally.

For a sheaf

\[
\mathcal{F},
\]

its cohomology groups are written

\[
\boxed{
H^i(X,\mathcal{F}).
}
\]

The symbol

\[
H^i
\]

is read ``H superscript \(i\).''

The zeroth group

\[
H^0(X,\mathcal{F})
\]

usually represents global sections.

Higher groups measure increasingly subtle global phenomena.

%%%%%%%%%%%%%%%%%%%%%%%%%%%%%%%%%%%%%%%%%%%%%%%%%%%%%%%%%%%%
\subsection{Algebraic Geometry as a Network}
\label{subsec:algebraic-geometry-network}
%%%%%%%%%%%%%%%%%%%%%%%%%%%%%%%%%%%%%%%%%%%%%%%%%%%%%%%%%%%%

The subject may be viewed schematically as

\begin{centereddiagram}
\begin{tikzpicture}[
node distance=11mm and 16mm,
every node/.style={
draw,
rounded corners,
align=center,
minimum width=3.5cm,
minimum height=8mm
}
]
\node (poly) {Polynomial\\Equations};
\node (var) [below=of poly] {Varieties};
\node (rings) [left=of var] {Coordinate\\Rings};
\node (proj) [right=of var] {Projective\\Geometry};
\node (scheme) [below=of var] {Schemes};
\node (local) [left=of scheme] {Local\\Algebra};
\node (sheaf) [right=of scheme] {Sheaves};
\node (arith) [below left=of scheme] {Arithmetic\\Geometry};
\node (bir) [below=of scheme] {Birational\\Geometry};
\node (complex) [below right=of scheme] {Complex\\Geometry};

\draw[-{Latex[length=3mm]}] (poly)--(var);
\draw[<->,>=Latex] (var)--(rings);
\draw[<->,>=Latex] (var)--(proj);
\draw[-{Latex[length=3mm]}] (var)--(scheme);
\draw[<->,>=Latex] (scheme)--(local);
\draw[<->,>=Latex] (scheme)--(sheaf);
\draw[-{Latex[length=3mm]}] (scheme)--(arith);
\draw[-{Latex[length=3mm]}] (scheme)--(bir);
\draw[-{Latex[length=3mm]}] (scheme)--(complex);
\end{tikzpicture}
\end{centereddiagram}

The classical theory begins with polynomial equations.

Modern algebraic geometry expands this viewpoint into a language capable of
studying arithmetic, topology, singularities, families, and higher-dimensional
spaces.

%%%%%%%%%%%%%%%%%%%%%%%%%%%%%%%%%%%%%%%%%%%%%%%%%%%%%%%%%%%%
\subsection{Common Errors}
\label{subsec:algebraic-geometry-common-errors}
%%%%%%%%%%%%%%%%%%%%%%%%%%%%%%%%%%%%%%%%%%%%%%%%%%%%%%%%%%%%

\subsubsection{Thinking Algebraic Geometry Is Only Geometry over \(\R\)}

Algebraic geometry frequently works over

\[
\C,
\qquad
\Q,
\qquad
\F_p,
\]

and arbitrary fields or commutative rings.

Real pictures represent only one part of the subject.

%%%%%%%%%%%%%%%%%%%%%%%%%%%%%%%%%%%%%%%%%%%%%%%%%%%%%%%%%%%%
\subsubsection{Confusing a Polynomial with Its Zero Set}

The polynomial

\[
f
\]

and the geometric set

\[
V(f)
\]

are different objects.

Different polynomials can define the same zero set.

For example,

\[
\boxed{
V(x)=V(x^2).
}
\]

%%%%%%%%%%%%%%%%%%%%%%%%%%%%%%%%%%%%%%%%%%%%%%%%%%%%%%%%%%%%
\subsubsection{Forgetting That \(V\) Reverses Inclusion}

If

\[
I\subseteq J,
\]

then

\[
\boxed{
V(J)\subseteq V(I).
}
\]

More equations produce fewer solutions.

%%%%%%%%%%%%%%%%%%%%%%%%%%%%%%%%%%%%%%%%%%%%%%%%%%%%%%%%%%%%
\subsubsection{Confusing Ideals with Radical Ideals}

Ordinary affine algebraic sets correspond naturally to radical ideals over
an algebraically closed field.

The ideals

\[
(x)
\]

and

\[
(x^2)
\]

are different even though they have the same classical zero set.

%%%%%%%%%%%%%%%%%%%%%%%%%%%%%%%%%%%%%%%%%%%%%%%%%%%%%%%%%%%%
\subsubsection{Assuming the Zariski Topology Resembles Euclidean Topology}

Zariski-open sets are generally much larger than Euclidean-open sets.

The topology is designed to reflect polynomial algebra rather than ordinary
distance.

%%%%%%%%%%%%%%%%%%%%%%%%%%%%%%%%%%%%%%%%%%%%%%%%%%%%%%%%%%%%
\subsubsection{Confusing Irreducible with Connected}

Irreducibility is stronger than connectedness.

An irreducible space cannot be decomposed into two proper closed subsets.

A connected space merely cannot be separated into two disjoint nonempty
open subsets.

%%%%%%%%%%%%%%%%%%%%%%%%%%%%%%%%%%%%%%%%%%%%%%%%%%%%%%%%%%%%
\subsubsection{Ignoring the Base Field}

The equation

\[
x^2+y^2+1=0
\]

has no real points but does have complex points.

Thus the notation

\[
V(f)
\]

must be interpreted relative to the chosen field.

%%%%%%%%%%%%%%%%%%%%%%%%%%%%%%%%%%%%%%%%%%%%%%%%%%%%%%%%%%%%
\subsubsection{Using Nonhomogeneous Equations in Projective Space}

An equation such as

\[
X^2+Y-1=0
\]

is not well-defined on projective equivalence classes.

Projective algebraic sets must be defined by homogeneous equations.

%%%%%%%%%%%%%%%%%%%%%%%%%%%%%%%%%%%%%%%%%%%%%%%%%%%%%%%%%%%%
\subsubsection{Assuming Every Algebraic Variety Is Smooth}

Polynomial equations can produce singularities.

Smoothness must be checked rather than assumed.

%%%%%%%%%%%%%%%%%%%%%%%%%%%%%%%%%%%%%%%%%%%%%%%%%%%%%%%%%%%%
\subsubsection{Confusing Classical Points with Scheme Points}

In scheme theory, points correspond to prime ideals, not only maximal ideals.

Generic points therefore appear naturally.

%%%%%%%%%%%%%%%%%%%%%%%%%%%%%%%%%%%%%%%%%%%%%%%%%%%%%%%%%%%%
\subsubsection{Assuming Schemes Are Merely Complicated Varieties}

Schemes genuinely retain more algebraic information.

They detect nilpotents, infinitesimal structure, arithmetic bases, and
families that classical point sets cannot fully represent.

%%%%%%%%%%%%%%%%%%%%%%%%%%%%%%%%%%%%%%%%%%%%%%%%%%%%%%%%%%%%
\subsection{Applications and Connections}
\label{subsec:algebraic-geometry-applications}
%%%%%%%%%%%%%%%%%%%%%%%%%%%%%%%%%%%%%%%%%%%%%%%%%%%%%%%%%%%%

\begin{applicationbox}
\textbf{Commutative Algebra.}
Ideals, prime ideals, maximal ideals, localization, dimension, and integral
extensions form much of the algebraic machinery underlying algebraic
geometry.

\medskip

\textbf{Number Theory.}
Polynomial equations over \(\Q\), number fields, and finite fields lead to
Diophantine and Arithmetic Geometry.

\medskip

\textbf{Projective Geometry.}
Projective space provides the natural global setting for homogeneous
polynomial equations and intersection theory.

\medskip

\textbf{Complex Geometry.}
Smooth complex algebraic varieties can be studied as complex manifolds,
connecting algebraic methods with analytic and differential-geometric ones.

\medskip

\textbf{Topology.}
Homology, cohomology, fundamental groups, and characteristic classes help
describe complex algebraic varieties.

\medskip

\textbf{Category Theory.}
Spectra reverse ring homomorphisms, fiber products describe geometric
constructions, and functorial methods organize much of the modern subject.

\medskip

\textbf{Cryptography.}
Elliptic curves and other algebraic varieties over finite fields are used in
public-key cryptography.

\medskip

\textbf{Coding Theory.}
Algebraic curves over finite fields produce important error-correcting codes.

\medskip

\textbf{Robotics and Kinematics.}
Polynomial constraints arise in linkage systems, inverse kinematics, and
configuration spaces.

\medskip

\textbf{Optimization.}
Polynomial optimization and semialgebraic geometry study optimization
problems defined by polynomial equations and inequalities.

\medskip

\textbf{Physics.}
Algebraic varieties, moduli spaces, complex geometry, and Calabi--Yau
varieties occur in several areas of mathematical physics and string theory.
\end{applicationbox}

%%%%%%%%%%%%%%%%%%%%%%%%%%%%%%%%%%%%%%%%%%%%%%%%%%%%%%%%%%%%
\subsection{Exercises}
\label{subsec:algebraic-geometry-exercises}
%%%%%%%%%%%%%%%%%%%%%%%%%%%%%%%%%%%%%%%%%%%%%%%%%%%%%%%%%%%%

\begin{exercise}[1]
Describe the real vanishing set

\[
V(y-x^2).
\]
\end{exercise}

\begin{exercise}[1]
Describe

\[
V(xy)
\subseteq\R^2.
\]
\end{exercise}

\begin{exercise}[1]
Find

\[
V(x,y)
\subseteq\A^2_k.
\]
\end{exercise}

\begin{exercise}[1]
Find the real vanishing set

\[
V(x^2+y^2).
\]
\end{exercise}

\begin{exercise}[2]
Compare the vanishing set of

\[
x^2+y^2
\]

over \(\R\) with its vanishing set over \(\C\).
\end{exercise}

\begin{exercise}[2]
Find

\[
V(x-y,x+y-4)
\subseteq\R^2.
\]
\end{exercise}

\begin{exercise}[2]
Show that

\[
V(I+J)=V(I)\cap V(J).
\]
\end{exercise}

\begin{exercise}[2]
Show that

\[
V(IJ)=V(I)\cup V(J).
\]
\end{exercise}

\begin{exercise}[2]
Explain why

\[
V(I)=V(\sqrt{I}).
\]
\end{exercise}

\begin{exercise}[2]
Show that

\[
V(x)=V(x^5).
\]
\end{exercise}

\begin{exercise}[2]
Determine

\[
\sqrt{(x^4)}
\]

in \(k[x]\).
\end{exercise}

\begin{exercise}[2]
Find the coordinate ring of

\[
V(y-x^3).
\]
\end{exercise}

\begin{exercise}[2]
Show that the coordinate ring of

\[
V(y-x^3)
\]

is isomorphic to \(k[x]\).
\end{exercise}

\begin{exercise}[2]
Explain why

\[
V(xy)
\]

is reducible.
\end{exercise}

\begin{exercise}[2]
Identify the irreducible components of

\[
V(xyz)
\subseteq\A^3.
\]
\end{exercise}

\begin{exercise}[2]
Describe the proper Zariski-closed subsets of

\[
\A^1_{\C}.
\]
\end{exercise}

\begin{exercise}[2]
For

\[
f(x)=x(x-1),
\]

describe the basic open set

\[
D(f)
\subseteq\A^1_{\R}.
\]
\end{exercise}

\begin{exercise}[2]
Determine whether the origin is singular on

\[
y^2=x^3.
\]
\end{exercise}

\begin{exercise}[2]
Determine whether the origin is singular on

\[
y-x^2=0.
\]
\end{exercise}

\begin{exercise}[2]
Find the tangent line to

\[
x^2+y^2=4
\]

at

\[
(0,2).
\]
\end{exercise}

\begin{exercise}[2]
Find the tangent line to

\[
y=x^2
\]

at

\[
(1,1).
\]
\end{exercise}

\begin{exercise}[3]
Compute the Zariski tangent space of

\[
V(y^2-x^3)
\]

at the origin.
\end{exercise}

\begin{exercise}[3]
Compare the dimension of the tangent space in the previous exercise with the
dimension of the curve.
\end{exercise}

\begin{exercise}[2]
Homogenize

\[
y=x^2+1.
\]
\end{exercise}

\begin{exercise}[2]
Homogenize

\[
x^3+y^2-1=0.
\]
\end{exercise}

\begin{exercise}[2]
Find the points at infinity of the projective closure of

\[
xy=1.
\]
\end{exercise}

\begin{exercise}[2]
Find the point at infinity of

\[
y=x^2.
\]
\end{exercise}

\begin{exercise}[2]
Determine the degree of each curve:

\[
X+Y-Z=0,
\]

\[
X^2+Y^2-Z^2=0,
\]

\[
Y^2Z-X^3+XZ^2=0.
\]
\end{exercise}

\begin{exercise}[2]
Use the genus formula to find the genus of a smooth plane conic.
\end{exercise}

\begin{exercise}[2]
Use the genus formula to find the genus of a smooth plane cubic.
\end{exercise}

\begin{exercise}[2]
Use the genus formula to find the genus of a smooth plane quintic.
\end{exercise}

\begin{exercise}[3]
Verify the rational parametrization

\[
x=\frac{1-t^2}{1+t^2},
\qquad
y=\frac{2t}{1+t^2}
\]

satisfies

\[
x^2+y^2=1.
\]
\end{exercise}

\begin{exercise}[3]
Explain why the rational parametrization of the circle misses or requires
special treatment of one point in its affine parameter.
\end{exercise}

\begin{exercise}[2]
Find the prime ideals of

\[
\Z
\]

that correspond to the ordinary prime numbers

\[
2,3,5,7.
\]
\end{exercise}

\begin{exercise}[2]
Explain the geometric interpretation of

\[
\Spec(\Z).
\]
\end{exercise}

\begin{exercise}[3]
Describe the difference between

\[
\Spec(k[x]/(x))
\]

and

\[
\Spec(k[x]/(x^2))
\]

from the scheme-theoretic viewpoint.
\end{exercise}

\begin{exercise}[3]
Explain why classical zero sets cannot detect the difference between
\((x)\) and \((x^2)\).
\end{exercise}

\begin{exercise}[3]
Explain what a generic point represents.
\end{exercise}

\begin{exercise}[3]
For the prime ideal

\[
(p)\subseteq\Z,
\]

determine its residue field.
\end{exercise}

\begin{exercise}[3]
Show that the residue field at

\[
(p)\in\Spec(\Z)
\]

is

\[
\F_p.
\]
\end{exercise}

\begin{exercise}[3]
Explain why a ring homomorphism

\[
R\to S
\]

induces a map

\[
\Spec(S)\to\Spec(R)
\]

rather than a map in the same direction.
\end{exercise}

\begin{exercise}[3]
Describe conceptually what the fiber of an arithmetic scheme over the prime
\(p\) represents.
\end{exercise}

\begin{exercise}[3]
Explain the analogy between coordinate charts on manifolds and affine open
subschemes.
\end{exercise}

\begin{exercise}[3]
Explain the role of the structure sheaf in an affine scheme.
\end{exercise}

\begin{exercise}[3]
What information does a local ring

\[
\mathcal{O}_{X,p}
\]

contain about the geometry near \(p\)?
\end{exercise}

\begin{exercise}[3]
Explain how

\[
\mathfrak m_p/\mathfrak m_p^2
\]

is related to tangent and cotangent spaces.
\end{exercise}

\begin{exercise}[3]
Explain why blowups are useful in the study of singularities.
\end{exercise}

\begin{exercise}[3]
State the geometric meaning of B\'ezout's theorem.
\end{exercise}

\begin{exercise}[3]
A line is tangent to a conic at a point. Explain why the intersection
multiplicity at that point can be \(2\).
\end{exercise}

\begin{exercise}[3]
Explain the meaning of birational equivalence.
\end{exercise}

\begin{exercise}[3]
Why do birationally equivalent irreducible varieties have isomorphic
function fields?
\end{exercise}

\begin{exercise}[3]
Explain what a moduli space attempts to classify.
\end{exercise}

\begin{exercise}[3]
What information does the \(j\)-invariant provide for elliptic curves over
an algebraically closed field of suitable characteristic?
\end{exercise}

\begin{exercise}[4]
Research the geometric interpretation of the Picard group.
\end{exercise}

\begin{exercise}[4]
Research the role of divisors in the Riemann--Roch theorem.
\end{exercise}

\begin{exercise}[4]
Investigate the relationship between algebraic curves over \(\C\) and compact
Riemann surfaces.
\end{exercise}

\begin{exercise}[4]
Research one example of a resolution of a plane-curve singularity.
\end{exercise}

\begin{exercise}[4]
Investigate how elliptic curves over finite fields are used in cryptography.
\end{exercise}

\begin{exercise}[4]
Research the role of schemes in modern number theory.
\end{exercise}

%%%%%%%%%%%%%%%%%%%%%%%%%%%%%%%%%%%%%%%%%%%%%%%%%%%%%%%%%%%%
\subsection{Conceptual Summary}
\label{subsec:algebraic-geometry-summary}
%%%%%%%%%%%%%%%%%%%%%%%%%%%%%%%%%%%%%%%%%%%%%%%%%%%%%%%%%%%%

Algebraic geometry begins with polynomial equations

\[
f_1=\cdots=f_m=0
\]

and studies their common solution sets as geometric spaces.

In affine space,

\[
\boxed{
V(I)
=
\{
p:
f(p)=0
\text{ for every }f\in I
\}.
}
\]

Conversely, a geometric set determines its vanishing ideal

\[
\boxed{
I(X)
=
\{
f:
f|_X=0
\}.
}
\]

Over an algebraically closed field, Hilbert's Nullstellensatz gives

\[
\boxed{
I(V(I))=\sqrt{I}.
}
\]

Thus radical ideals correspond to classical affine algebraic sets.

The Zariski topology declares algebraic sets to be closed.

Irreducible algebraic sets correspond to prime ideals.

The coordinate ring

\[
\boxed{
k[X]
=
k[x_1,\ldots,x_n]/I(X)
}
\]

encodes polynomial functions on an affine variety.

Projective geometry replaces affine coordinates with homogeneous coordinates

\[
[X_0:\cdots:X_n]
\]

and provides natural points at infinity.

Smoothness can be studied using Jacobian matrices and tangent spaces.

For an algebraic set defined by

\[
f_1,\ldots,f_m,
\]

the Zariski tangent space is determined by the linearized equations

\[
\boxed{
J(p)v=0.
}
\]

Algebraic curves possess invariants such as degree and genus.

Rational functions form the function field

\[
\boxed{
k(X).
}
\]

Birational geometry studies varieties with the same rational function field.

Modern algebraic geometry replaces classical varieties with schemes.

For a commutative ring \(R\),

\[
\boxed{
\Spec(R)
=
\{
\text{prime ideals of }R
\}.
}
\]

A scheme combines this topological space with a structure sheaf containing
local algebraic information.

The fundamental philosophy is therefore

\[
\boxed{
\text{geometry}
\longleftrightarrow
\text{commutative algebra}.
}
\]

This correspondence makes algebraic geometry a meeting point of algebra,
number theory, topology, geometry, and analysis.

%%%%%%%%%%%%%%%%%%%%%%%%%%%%%%%%%%%%%%%%%%%%%%%%%%%%%%%%%%%%
\subsection{Section Connections}
\label{subsec:algebraic-geometry-connections}
%%%%%%%%%%%%%%%%%%%%%%%%%%%%%%%%%%%%%%%%%%%%%%%%%%%%%%%%%%%%

\chapterconnection{
Algebraic Geometry transforms polynomial algebra into geometry. Its
coordinate rings, ideals, localizations, prime ideals, and dimensions rely
directly on Commutative Algebra. Projective Geometry supplies the natural
global setting for homogeneous polynomial equations. Diophantine Geometry
and Arithmetic Geometry ask arithmetic questions about algebraic varieties
over fields such as \(\Q\), \(\Q_p\), and \(\F_p\). Complex algebraic
varieties connect the subject with Complex Geometry, Differential Geometry,
and topology. Schemes, sheaves, fiber products, and functorial constructions
connect algebraic geometry with Category Theory and Homological Algebra.
The next section, Complex Geometry, develops geometric spaces whose local
coordinates and transition maps are complex analytic.
}

%%%%%%%%%%%%%%%%%%%%%%%%%%%%%%%%%%%%%%%%%%%%%%%%%%%%%%%%%%%%
% SECTION SOURCE TRACKING
%%%%%%%%%%%%%%%%%%%%%%%%%%%%%%%%%%%%%%%%%%%%%%%%%%%%%%%%%%%%

%------------------------------------------------------------
% 3.8 Algebraic Geometry
%------------------------------------------------------------
%
% Source basis:
%   - Standard classical and introductory modern algebraic
%     geometry.
%
% Used for:
%   - affine space and polynomial zero sets;
%   - vanishing ideals and radical ideals;
%   - Nullstellensatz;
%   - Zariski topology;
%   - irreducibility and prime ideals;
%   - coordinate rings;
%   - regular functions and morphisms;
%   - dimension and hypersurfaces;
%   - algebraic curves and surfaces;
%   - singularities and Zariski tangent spaces;
%   - projective varieties and homogenization;
%   - Bezout's theorem and intersection multiplicity;
%   - rational functions and birational equivalence;
%   - genus and divisors;
%   - Riemann--Roch overview;
%   - local rings;
%   - spectra and affine schemes;
%   - structure sheaves;
%   - residue fields and generic points;
%   - morphisms, fibers, and base change;
%   - Kahler differentials;
%   - resolution and blowups;
%   - moduli and birational geometry;
%   - sheaf cohomology overview.
%
% Notes:
%   - Ideal theory, localization, prime spectra, and Krull
%     dimension are developed canonically in Commutative
%     Algebra.
%   - Projective coordinates and incidence geometry are
%     developed canonically in Projective Geometry.
%   - Arithmetic questions concerning rational and integral
%     points belong primarily to Diophantine Geometry and
%     Arithmetic Geometry.
%   - Complex manifolds, Hermitian geometry, and Kahler
%     geometry are developed in Complex Geometry.
%   - Sheaf cohomology and derived methods are expanded in
%     Homological Algebra and Algebraic Topology.
%   - Scheme theory is introduced here because it is part of
%     the modern language of algebraic geometry, but advanced
%     scheme theory is beyond the introductory development.
%
%%%%%%%%%%%%%%%%%%%%%%%%%%%%%%%%%%%%%%%%%%%%%%%%%%%%%%%%%%%%